% Content of the blueprint, shared by the web and print versions.
%
% The chapter files in chapters/ are generated: every node's \lean, \leanok and \uses
% are computed from the compiled Lean project (the constants occurring in each
% declaration's type and value, followed through declarations that are not nodes of
% this blueprint). Statement dependencies go in the statement's \uses, proof
% dependencies in the proof's \uses.

\chapter*{How to read this blueprint}

This blueprint accompanies the Lean~4 formalization
\href{https://github.com/SamPease/GSpinGGP}{GSpinGGP} (work in progress) of the dissertation
\cite{Thesis}: S.~Pease, \emph{The Local Gan--Gross--Prasad Conjecture for General Spin Groups},
Rutgers University--Newark, 2026, \href{https://doi.org/10.7282/t3-qb3s-s249}{doi:10.7282/t3-qb3s-s249}.
It is cited as ``Thesis'', with the numbers of the printed dissertation. The proof adapts
Beuzart-Plessis's local trace formula for unitary groups \cite{BP} and the work of Waldspurger
\cite{Wal10} and Luo \cite{Luo} on special orthogonal groups.

Every node is a Lean declaration, and every statement is formalized. The main theorem
(Theorem~\ref{local_ggp}) is proved in Lean from three kinds of inputs: results whose Lean proof is
still missing (\texttt{sorry}), results from the literature that are assumed (Lean \texttt{axiom}s),
and definitions. Results are attributed as follows.
\begin{itemize}
  \item \textbf{Proved}: the Lean proof of this step contains no \texttt{sorry}. It may still rely,
    through the results it uses, on unproved steps or on the literature.
  \item \textbf{Not yet proved}: the Lean proof is \texttt{sorry}. Its proof paragraph gives the
    intended argument.
  \item \textbf{Assumed from the literature}: a ``Literature'' environment, with an
    ``Assumed from'' line giving the source. These are Lean \texttt{axiom}s; they have no proof
    paragraph, and are never proved in this project. Several are adapted to general spin groups from
    statements for unitary or special orthogonal groups.
  \item \textbf{Definitions}: real Lean definitions (never axiomatized).
\end{itemize}
Where a Lean statement deliberately differs from the dissertation, its node says so in one line.

\paragraph{Colours in the dependency graph.}
Boxes are definitions, ellipses are results. A green border means the statement is formalized,
which is the case for every node. The interior colour gives the status of the proof:
\begin{itemize}
  \item \emph{light green fill}: proved (the Lean proof of this step has no \texttt{sorry});
  \item \emph{dark green fill}: proved, and so is everything it depends on (no \texttt{sorry} and no
    literature assumption among its ancestors);
  \item \emph{blue fill}: not yet proved (\texttt{sorry}); the statement and everything it uses are
    formalized, so the proof is ready to be written;
  \item \emph{no fill (white)}, with a label starting with \texttt{LIT-}: assumed from the literature
    (Lean \texttt{axiom});
  \item \emph{light green box}: a definition.
\end{itemize}
Because every result on the path of the main theorem rests on literature assumptions, none of them is
dark green: the dark green nodes are the self-contained results. The graph's own legend uses the
generic wording of the blueprint software (``ready to be formalized'' for the blue fill).

\paragraph{Status} (October 2026). The project builds with 29 declarations containing
\texttt{sorry}; the main theorem depends on 18 of them and on 28 of the 31 literature assumptions.
Classical facts about quadratic forms (Witt cancellation and decomposition, Cartan--Dieudonn\'e,
the classification over $p$-adic fields) are proved from the Tau Ceti library.
The 11 unproved statements not on its path are collected in Chapter~\ref{chap:offpath}. To see exactly
what the main theorem rests on, run \texttt{\#print axioms GSpinGGP.local\_ggp}.

% Generated from the Lean dependency data; see the comment at the top of content.tex.
\chapter{General spin groups and GGP triples}
\label{chap:gspin}
This chapter follows Thesis Ch.~5. Throughout, $F$ is a $p$-adic field (a finite extension of
$\mathbb{Q}_p$) and $\psi$ a nontrivial continuous additive character of $F$. Quadratic spaces are
nondegenerate and finite-dimensional; $b$ is the bilinear form with $q(v)=b(v,v)$.

\section{Quadratic spaces and the general spin group}
\begin{definition}[The general spin group]
\label{GSpin}
\lean{GSpinGGP.GSpin}
\leanok
For a quadratic space $(V,q)$ with Clifford algebra $C(V)$ and even part $C^+(V)$, the general spin
group is
\[ \GSpin(V) = \{ g \in C^+(V)^\times : g V g^{-1} = V \} \]
(Thesis \S5.1.2). Conjugation gives a surjection $P:\GSpin(V)\to \SO(V)$ on $\Fbar$-points with
kernel the scalars $\mathbb{G}_m$. (This is not the spin group of unit spinor norm.)
\end{definition}

\begin{definition}[The Lie algebra $\mathfrak{gspin}(V)$]
\label{gspin}
\lean{GSpinGGP.gspin}
\leanok
\uses{GSpin}
$\gspin(V) = \{ X \in C^+(V) : [X,v] \in V \text{ for all } v\in V\}$, a Lie subalgebra of
$C(V)$ for the commutator bracket (Thesis Prop.~5.1.3.1); $\gspin(V)=\mathfrak{so}(V)\oplus F$.
\end{definition}

\begin{theorem}[Witt decomposition]
\label{witt_decomposition}
\lean{GSpinGGP.witt_decomposition}
\leanok
A quadratic space is $V \cong H^{w}\perp V_{\mathrm{an}}$ with $H$ the hyperbolic plane and
$V_{\mathrm{an}}$ anisotropic, where $w=w(V)$ is the Witt index (the maximal dimension of a totally
isotropic subspace). In particular $\dim V_{\mathrm{an}} = \dim V - 2w(V)$, so $\GSpin(V)$ is
quasi-split iff $\dim V - 2w(V)\le 2$.
\end{theorem}
\begin{proof}
\leanok
Proved from Tau Ceti's Witt decomposition of a regular form (\cite{Lam}, Ch.~I \S4).
\end{proof}

\begin{theorem}[Witt cancellation]
\label{witt_cancellation}
\lean{GSpinGGP.witt_cancellation}
\leanok
Over a field of characteristic $\neq 2$, if $U\perp X \cong U\perp Y$ then $X\cong Y$
(all forms finite-dimensional, possibly degenerate).
\end{theorem}
\begin{proof}
\leanok
Split off the radicals and apply Tau Ceti's Witt cancellation for a regular summand (\cite{Lam}, Ch.~I \S4).
\end{proof}

\begin{theorem}[Quadratic forms over a $p$-adic field]
\label{padic_quadratic_classes}
\lean{GSpinGGP.padic_quadratic_classes}
\leanok
Let $Q$ be a nondegenerate quadratic form on $F^m$. The isometry classes of nondegenerate forms
on $F^m$ with the same discriminant as $Q$ are finite in number, $n\le 2$ of them, and admit a system of
representatives containing $Q$; $n=2$ if $m\ge 3$, $n=1$ if $m\le 1$, and for $m=2$, $n=1$ iff
$-\det Q$ is a square.
\end{theorem}
\begin{proof}
\leanok
Proved from Tau Ceti's classification of regular forms over a nonarchimedean local field by rank, discriminant and Hasse invariant, and its list of the realizable invariants (\cite{Lam}, Ch.~VI \S2; O'Meara \S63), through the fact that a $p$-adic field is a nonarchimedean local field.
\end{proof}

\begin{theorem}[Cartan--Dieudonn\'e, $\SO$ form]
\label{cartan_dieudonne_even}
\lean{GSpinGGP.cartan_dieudonne_even}
\leanok
In a nondegenerate quadratic space over a field of characteristic $\neq 2$, every isometry of
determinant $1$ is a product $\prod_i s_{u_i}s_{w_i}$ of pairs of reflections in anisotropic vectors.
\end{theorem}
\begin{proof}
\leanok
Proved from Tau Ceti's factorization of an element of $\SO$ into an even number of reflections (\cite{Artin}, Ch.~III; \cite{Lam}, Ch.~I), grouped in pairs.
\end{proof}

\begin{literature}[Centralizers of semisimple elements of $\GSpin$]
\label{LIT-steinberg_centralizer_connected}
\lean{GSpinGGP.steinberg_centralizer_connected}
\leanok
\uses{GSpin}
Let $x\in\GSpin(V)$ be semisimple, $\bar x=P(x)$, and $W_{\pm1}=\ker(\bar x\mp 1)$. Then
$Z_{\GSpin(V)}(x)=P^{-1}(Z_{\SO(V)}(\bar x)^0)$; explicitly, $g$ commutes with $x$ iff $P(g)$
commutes with $\bar x$ and $\det P(g)|_{W_{+1}} = \det P(g)|_{W_{-1}} = 1$. (Characteristic $0$.)

\assumedfrom{\cite{Steinberg}, Thm.~8.1 (the derived group $\Spin(V)$ is simply connected)}
\end{literature}

\begin{literature}[Regular nilpotent orbits and quasi-splitness]
\label{LIT-nilReg_nonempty_iff_quasiSplit}
\lean{GSpinGGP.nilReg_nonempty_iff_quasiSplit}
\leanok
\uses{GSpin, gspin}
$\Nilreg(\gspin(V))\neq\emptyset$ iff $\GSpin(V)$ is quasi-split, i.e.\ $\dim V-2w(V)\le 2$.

\assumedfrom{\cite{Wal01}, Ch.~I; \cite{Wal10}, \S7.1; \cite{Luo}, \S2.1; Thesis \S5.1.4}
\end{literature}

\section{GGP triples}
\begin{definition}[Admissible pairs]
\label{GGPData}
\lean{GSpinGGP.GGPData}
\leanok
A GGP datum is a quadratic space $W$ on $F^m$ together with $r\ge 0$ and $\nu\in F^\times$. Put
$Z$ the split space of dimension $2r+1$ with basis $z_{-r},\dots,z_r$, $b(z_i,z_j)=\nu\delta_{i,-j}$,
and $V = W\oplus^\perp Z$ (Thesis \S5.2).
\end{definition}

\begin{definition}[The group $G$]
\label{G}
\lean{GSpinGGP.GGPData.G}
\leanok
\uses{GGPData, GSpin}
$G = \GSpin(W)\times\GSpin(V)$ (Thesis \S5.2), as a subgroup of $(C(W)\times C(V))^\times$.
Its Lie algebra is $\mathfrak g = \gspin(W)\oplus\gspin(V)$, and $A_G$ denotes its split centre.
\end{definition}

\begin{definition}[The unipotent subgroup $N$]
\label{N}
\lean{GSpinGGP.GGPData.N}
\leanok
\uses{GGPData, GSpin}
$N\subset \GSpin(V)$ is the unipotent radical of the stabilizer $P_V$ in $\GSpin(V)$ of the isotropic
flag $\langle z_r\rangle\subset\langle z_r,z_{r-1}\rangle\subset\dots\subset\langle z_r,\dots,z_1\rangle$
(Thesis \S5.2).
\end{definition}

\begin{definition}[The subgroup $H$]
\label{H}
\lean{GSpinGGP.GGPData.H}
\leanok
\uses{G, GGPData, GSpin, N}
$H = \GSpin(W)\ltimes N\subset G$: the pairs $(a, an)$ with $a\in\GSpin(W)$, $n\in N$
(Thesis \S5.2). $Z_{G,H}=A_{G,H}$ is the diagonal central $\mathbb{G}_m$.
\end{definition}

\begin{definition}[The character $\xi$]
\label{xi}
\lean{GSpinGGP.GGPData.xi}
\leanok
\uses{G, GGPData, GSpin, H}
$\xi(h) = \psi(\lambda(h))$, where $\lambda(an) = \sum_{i=0}^{r-1} b(z_{-i-1}, n z_i)$ is a
homomorphism $H\to\mathbb{G}_a$ (Thesis \S5.2).
\end{definition}

\begin{definition}[Quasi-split triples]
\label{IsQuasiSplit}
\lean{GSpinGGP.GGPData.IsQuasiSplit}
\leanok
\uses{GGPData}
$G$ is quasi-split iff both $W$ and $V$ are, i.e.\ $\dim U-2w(U)\le 2$ for $U=W,V$ (Thesis \S6.4).
\end{definition}

\begin{definition}[Central characters]
\label{CentralCharOn}
\lean{GSpinGGP.CentralCharOn}
\leanok
\uses{G, GGPData, GSpin}
A continuous unitary character $\chi: A_G(F)\to\mathbb{S}^1$ trivial on $A_{G,H}(F)$
(Thesis \S7.1). It is fixed throughout the trace formula.
\end{definition}

\begin{definition}[Whittaker data]
\label{WhittakerDatum}
\lean{GSpinGGP.WhittakerDatum}
\leanok
\uses{G, GGPData, GSpin}
A Whittaker datum of $G$ (Thesis \S4.4): a Borel subgroup $B$ (stabilizer of full maximal isotropic
flags of $W$ and $V$) with unipotent radical $U$, and a generic continuous unitary character $\chi$
of $U(F)$, i.e.\ its stabilizer in $B(F)$ is $Z_G(F)U(F)$. (The Thesis takes $G(F)$-conjugacy classes
of such pairs.)
\end{definition}

\begin{literature}[Luo's orbit $\mathcal O_\nu$]
\label{LIT-exists_regularNil_inOrbitNu}
\lean{GSpinGGP.exists_regularNil_inOrbitNu}
\leanok
\uses{GGPData, IsQuasiSplit, gspin}
Let $G$ be quasi-split. If $\dim V$ is even and $\ge 4$, some regular nilpotent $X\in\gspin(V)$
has image in $\mathfrak{so}(V)$ in the orbit $\mathcal O_\nu$ (it kills an anisotropic line $Fd_0$ and is
one Jordan block on $U_1=d_0^\perp\cong\langle\nu\rangle\oplus H^k$). If $\dim W$ is even and $\ge 4$,
some regular nilpotent $X\in\gspin(W)$ has image in $\mathcal O_{-\nu}$.

\assumedfrom{\cite{Wal10}, \S7.1; \cite{Luo}, \S2.1, \S7.1; \cite{Wal01}, Ch.~I; Thesis \S5.1.4}
\end{literature}

\begin{literature}[Germs vanish off quasi-split centralizers]
\label{LIT-nilReg_centralizer_nonempty_imp}
\lean{GSpinGGP.nilReg_centralizer_nonempty_imp}
\leanok
\uses{G, GGPData, GSpin, gspin}
For semisimple $x\in\GSpin(W)(F)$ embedded diagonally in $G$, with $W'_x=\ker(P(x)-1)$ and
$V'_x = W'_x\oplus Z$: if $\Nilreg(\mathfrak g_x)\ne\emptyset$ then $W'_x$ and $V'_x$ are quasi-split.

\assumedfrom{\cite{Wal01}, Ch.~I applied to the factors of $G_x$; \cite{BP}, Prop.~4.5.1(i)}
\end{literature}

\begin{literature}[Regular nilpotent orbits of GGP centralizers]
\label{LIT-nilReg_centralizer_nonempty_iff}
\lean{GSpinGGP.nilReg_centralizer_nonempty_iff}
\leanok
\uses{G, GGPData, GSpin, gspin}
For semisimple $x\in\GSpin(W)(F)$ whose moved part $G''_x$ is a torus (true for every $x$ contributing
to $\Gamma(G,H)$): $\Nilreg(\mathfrak g_x)\neq\emptyset$ iff $W'_x$ and $V'_x=W'_x\oplus Z$ are
quasi-split.

\assumedfrom{\cite{Wal01}, Ch.~I; \cite{Wal10}, \S7.1}
\end{literature}

\begin{literature}[Multiplicity at most one]
\label{LIT-emoryTakeda_mult_le_one}
\lean{GSpinGGP.emoryTakeda_mult_le_one}
\leanok
\uses{G, GGPData, GSpin, H, IsTempered, UnitaryRep, xi}
For every GGP triple and every irreducible tempered unitary representation $\pi$ of $G(F)$,
$m(\pi)=\dim\Hom_{H(F)}(\pi^\infty,\xi)\le 1$ (Thesis Thm.~5.5.1).

\assumedfrom{\cite{ET}, Thm.~1.3 ($r=0$), with the reduction of \cite{BP}, \S6.3 and \cite{GGP}, \S15; directly \cite{Yan}}
\end{literature}


% Generated from the Lean dependency data; see the comment at the top of content.tex.
\chapter{Representations, characters and quasi-characters}
\label{chap:reps}
Background of Thesis Ch.~2--4: tempered representations, the multiplicity $m(\pi)$,
Harish-Chandra characters, quasi-characters and their germs.

\section{Representations and the multiplicity}
\begin{definition}[Unitary representations]
\label{UnitaryRep}
\lean{GSpinGGP.UnitaryRep}
\leanok
A strongly continuous unitary representation of a topological group on a complex Hilbert space
(Thesis \S3.1).
\end{definition}

\begin{definition}[Tempered representations]
\label{IsTempered}
\lean{GSpinGGP.IsTempered}
\leanok
\uses{UnitaryRep}
$\pi$ is irreducible and tempered: irreducible and weakly contained in the left regular representation
of $G$ (Thesis Ch.~1, \S3.1.2). The Thesis uses the growth condition on matrix coefficients and notes the
equivalence (Cowling--Haagerup--Howe); the weak-containment form is used here because it stays valid
when the centre of $G$ is noncompact.
\end{definition}

\begin{definition}[The multiplicity $m(\pi)$]
\label{mult}
\lean{GSpinGGP.mult}
\leanok
\uses{G, GGPData, H, UnitaryRep, xi}
For a unitary representation $\pi$ of $G(F)$, \[ m(\pi)=\dim\Hom_{H(F)}(\pi^\infty,\xi)\in\mathbb N\cup\{\infty\}, \]
where $\pi^\infty$ is the space of smooth vectors (Thesis \S5.5).
\end{definition}

\begin{literature}[Schur's lemma]
\label{LIT-schur_central}
\lean{GSpinGGP.schur_central}
\leanok
\uses{UnitaryRep}
In an irreducible unitary representation of a topological group, the centre acts by a character.

\assumedfrom{\cite{Dixmier}; Thesis \S3.1.2}
\end{literature}

\begin{literature}[$G(F)$ is unimodular]
\label{LIT-haar_G_isMulRightInvariant}
\lean{GSpinGGP.LitTF.haar_G_isMulRightInvariant}
\leanok
\uses{G, GGPData, GSpin}
The Haar measure of $G(F)=\GSpin(W)(F)\times\GSpin(V)(F)$ (and of every inner form) is right invariant.

\assumedfrom{Bourbaki, \emph{Groupes et alg\`ebres de Lie} III \S3; \cite{BP}, \S1.6; Thesis \S2.4}
\end{literature}

\begin{literature}[Weil's quotient-measure formula]
\label{LIT-weil_quotient}
\lean{GSpinGGP.weil_quotient}
\leanok
Let $S\hookrightarrow\Gamma$ be a closed embedding of unimodular locally compact, Hausdorff, second countable
groups with Haar measures $\mu,\nu$. There is a unique measure $m$ on $S\backslash\Gamma$ with
$\int_\Gamma\varphi\,d\nu=\int_{S\backslash\Gamma}\int_S\varphi(sq)\,d\mu(s)\,dm(q)$ for all
$\varphi\in C_c(\Gamma)$. (Used for $H(F)/A_{G,H}(F)\subset G(F)/A_G(F)$, Thesis \S7.1.)

\assumedfrom{Bourbaki, \emph{Int\'egration} VII \S2; \cite{Folland}, Thm.~2.49}
\end{literature}

\begin{literature}[Integration over $G=PK$]
\label{LIT-integral_eq_mul_integral_compact_integral}
\lean{GSpinGGP.integral_eq_mul_integral_compact_integral}
\leanok
Let $G$ be a locally compact, Hausdorff, second countable group, $P$ a closed subgroup and $K$ a compact
subgroup with $G=PK$. For Haar measures $dx,dp,dk$ there is $c>0$ with
$\int_G f(x)\,dx = c\int_K\int_P f(pk)\,dp\,dk$ for continuous compactly supported $f$.

\assumedfrom{\cite{Knapp}, Thm.~8.32; Bourbaki, \emph{Int\'egration} VII \S2; Casselman's notes \S3.1}
\end{literature}

\section{Characters and quasi-characters}
\begin{definition}[The Weyl discriminant]
\label{weylDisc}
\lean{GSpinGGP.weylDisc}
\leanok
$D^G(x)=|\det(1-\Ad x)|_{\mathfrak g/\mathfrak g_x}|_F$ (Thesis \S2.1), computed as the absolute value of the
trailing coefficient of the characteristic polynomial of $1-\Ad x$ on $\mathfrak g$.
\end{definition}

\begin{literature}[The regular set]
\label{LIT-hc_regularSet_open_conull}
\lean{GSpinGGP.hc_regularSet_open_conull}
\leanok
\uses{G, GGPData, GSpin, gspin}
The regular semisimple set $\Greg(F)$ is open in $G(F)$ and its complement has Haar measure zero.

\assumedfrom{\cite{HC70}; \cite{BP}, \S1; Thesis \S2.1}
\end{literature}

\begin{literature}[Harish-Chandra's regularity theorem]
\label{LIT-hc_character_exists}
\lean{GSpinGGP.hc_character_exists}
\leanok
\uses{G, GGPData, GSpin, IsTempered, UnitaryRep, gspin}
For an irreducible tempered $\pi$, the distribution $f\mapsto\operatorname{tr}\pi(f)$ on $C_c^\infty(G(F))$
is given by a locally integrable function which is locally constant on $\Greg(F)$.

\assumedfrom{Bernstein (admissibility); \cite{HC99}; \cite{BP}, \S2.1; Thesis Ch.~1}
\end{literature}

\begin{definition}[The character $\theta_\pi$]
\label{charG}
\lean{GSpinGGP.charG}
\leanok
\uses{G, GGPData, GSpin, UnitaryRep, gspin, mult}
The Harish-Chandra character $\theta_\pi$ of $\pi$, the representative of $f\mapsto\operatorname{tr}\pi(f)$
that is locally constant on $\Greg(F)$ (Thesis Ch.~1).
\end{definition}

\begin{definition}[Fourier transforms of nilpotent orbits]
\label{jHat}
\lean{GSpinGGP.jHat}
\leanok
For a nilpotent orbit $\mathcal O\subset\mathfrak g(F)$, $\hat{j}(\mathcal O,\cdot)$ is the function
representing $f\mapsto\mu_{\mathcal O}(\hat f)$, the Fourier transform being normalized by $b$ and $\psi$
(\cite{BP} \S1.8, \S4.1; Thesis \S4.1.1). It is unique on $\mathfrak g_{\mathrm{reg}}$.
\end{definition}

\begin{definition}[Quasi-characters]
\label{IsQuasiCharacter}
\lean{GSpinGGP.IsQuasiCharacter}
\leanok
\uses{jHat}
$\theta: G(F)\to\mathbb C$ is a quasi-character (Thesis \S4.1.1, \cite{BP} \S4.1) if it is conjugation
invariant on $\Greg(F)$ and at every semisimple $x$ there are coefficients $c_{\theta,\mathcal O}(x)$,
$\mathcal O\in\Nil(\mathfrak g_x)$, with
$\theta(xe^X)=\sum_{\mathcal O}c_{\theta,\mathcal O}(x)\,\hat j(\mathcal O,X)$ for $X\in\mathfrak g_x(F)$
regular and near $0$.
\end{definition}

\begin{definition}[Germs]
\label{germ}
\lean{GSpinGGP.germ}
\leanok
\uses{jHat}
For a quasi-character $\theta$, semisimple $x$ and $\mathcal O\in\Nilreg(\mathfrak g_x)$, the germ
$c_{\theta,\mathcal O}(x)$ is the coefficient of $\hat j(\mathcal O,\cdot)$ in the local expansion of $\theta$
at $x$ (Thesis \S4.2, \cite{BP} \S4.5); it is $0$ if $\mathcal O$ is not regular or there is no expansion.
\end{definition}

\begin{literature}[Uniqueness of germs]
\label{LIT-germ_unique}
\lean{GSpinGGP.germ_unique}
\leanok
\uses{G, GGPData, gspin, jHat}
At semisimple $x$, the coefficients of a local expansion $\theta(xe^X)=\sum_{\mathcal O}c_{\mathcal O}\hat j(\mathcal O,X)$
are unique on the regular orbits $\mathcal O\in\Nilreg(\mathfrak g_x)$ (linear independence of the
$\hat j(\mathcal O,\cdot)$ near $0$).

\assumedfrom{\cite{HC99}; \cite{BP}, \S4.5; Thesis \S4.2}
\end{literature}

\begin{literature}[Local character expansion]
\label{LIT-hc_character_isQuasiCharacter}
\lean{GSpinGGP.hc_character_isQuasiCharacter}
\leanok
\uses{G, GGPData, GSpin, IsQuasiCharacter, IsTempered, UnitaryRep, gspin}
For irreducible tempered $\pi$, $\theta_\pi$ is a quasi-character.

\assumedfrom{\cite{HC99}; \cite{Wal10}, \S4; \cite{BP}, Prop.~4.1.1(i); Thesis Prop.~4.1.1.1(i)}
\end{literature}

\begin{literature}[Rodier's theorem]
\label{LIT-rodier_germ_one}
\lean{GSpinGGP.rodier_germ_one}
\leanok
\uses{G, GGPData, GSpin, IsQuasiSplit, IsTempered, UnitaryRep, WhittakerDatum, germ, gspin}
Let $G$ be quasi-split. For each $\mathcal O\in\Nilreg(\mathfrak g)$ there is a Whittaker datum $(U,\chi)$
such that for every irreducible tempered $\pi$, $c_{\theta_\pi,\mathcal O}(1)=1$ if $\pi$ has a
$(U,\chi)$-Whittaker model and $0$ otherwise (Thesis Prop.~4.4.1(i)).

\assumedfrom{\cite{Rodier}; \cite{MW87}; \cite{BP}, Prop.~4.8.1(i)}
\end{literature}

\begin{literature}[Self-dual Haar measures]
\label{LIT-exists_isSelfDual}
\lean{GSpinGGP.exists_isSelfDual}
\leanok
For a finite-dimensional $F$-vector space $V$ with a nondegenerate bilinear form $B$ there is a Haar measure
$\mu$ on $V$ for which Fourier inversion holds: $f(Z)=\int\hat f(Y)\psi(-B(Z,Y))\,d\mu(Y)$ with
$\hat f(Y)=\int f(X)\psi(B(X,Y))\,d\mu(X)$, for all $f\in C_c^\infty(V)$.

\assumedfrom{Tate's thesis; \cite{Weil}, Ch.~II \S5, Ch.~VII \S2; \cite{BP}, \S1.6; Thesis \S2.4}
\end{literature}


% Generated from the Lean dependency data; see the comment at the top of content.tex.
\chapter{Inner forms and the main theorem}
\label{chap:mainthm}
Thesis \S6.1--6.4. The Langlands parameter is not modelled: the main theorem is proved for every
family of packets satisfying the stability, transfer and genericity hypotheses, which is all the proof
uses.

\section{Inner forms, stability and packets}
\begin{definition}[Inner twists]
\label{InnerTwist}
\lean{GSpinGGP.InnerTwist}
\leanok
An inner twist of $\GSpin(W)$: a quadratic space $W'$ on $F^m$ of the same dimension and discriminant
with a fixed $\Fbar$-isometry $\varphi:W_{\Fbar}\simeq W'_{\Fbar}$ such that
$\varphi^{-1}\circ\sigma(\varphi)\in\SO(W)(\Fbar)$ for all $\sigma$. It gives
$\psi_\alpha=C(\varphi)\times C(\varphi\oplus\mathrm{id}_Z):G_{\Fbar}\simeq G_{\alpha,\Fbar}$.
This differs from Thesis \S6.2: the inner forms are indexed by $H^1(F,\SO(W))$, not $H^1(F,\mathrm{PGSpin}(W))$.
\end{definition}

\begin{definition}[Stable conjugacy across inner forms]
\label{StablyConjG}
\lean{GSpinGGP.GGPData.StablyConjG}
\leanok
\uses{G, GGPData, GSpin, InnerTwist}
$x\in G(F)$ and $y\in G_\alpha(F)$ are stably conjugate if $y=\psi_\alpha(g x g^{-1})$ for some
$g\in G(\Fbar)$, componentwise.
\end{definition}

\begin{definition}[The Kottwitz sign]
\label{kottwitzSign}
\lean{GSpinGGP.kottwitzSign}
\leanok
\uses{GGPData, InnerTwist}
$e(G_\alpha) = (-1)^{\mathrm{rk}_F G-\mathrm{rk}_F G_\alpha}=(-1)^{w(W)+w(W_\alpha)+w(V)+w(V_\alpha)}$
(the $p$-adic formula; that it equals Kottwitz's cohomological sign is \cite{Kot83}).
\end{definition}

\begin{definition}[Systems of inner forms]
\label{InnerFormSystem}
\lean{GSpinGGP.InnerFormSystem}
\leanok
\uses{GGPData, InnerTwist}
An index set $\{\alpha\}$ with inner twists $(W_\alpha,\varphi_\alpha)$, $\alpha=1$ the trivial one, such that every
nondegenerate form on $F^m$ with the discriminant of $W$ is isometric to exactly one $W_\alpha$; so the index
set is $H^1(F,\SO(W))=H^1(F,H/Z_{G,H})$. $(G_\alpha,H_\alpha,\xi_\alpha)$ is the GGP triple of
$(W_\alpha,V_\alpha=W_\alpha\oplus Z)$. This differs from Thesis \S6.2: the inner forms are indexed by
$H^1(F,\SO(W))$, not $H^1(F,\mathrm{PGSpin}(W))$.
\end{definition}

\begin{definition}[Stable functions]
\label{IsStableFn}
\lean{GSpinGGP.IsStableFn}
\leanok
\uses{G, GGPData, GSpin, InnerTwist, gspin}
(STAB) $\theta$ is constant on stable conjugacy classes of regular semisimple elements (Thesis \S6.1).
\end{definition}

\begin{definition}[Tempered packets on the inner forms]
\label{TemperedPacketFamily}
\lean{GSpinGGP.TemperedPacketFamily}
\leanok
\uses{G, GGPData, GSpin, InnerFormSystem, InnerTwist, IsStableFn, IsTempered, StablyConjG, UnitaryRep, WhittakerDatum, charG, gspin, kottwitzSign}
Finite families $\Pi_\alpha$ of pairwise inequivalent irreducible tempered representations of $G_\alpha(F)$
with central characters trivial on $Z_{G,H}(F)$, such that, with $\theta_\alpha=\sum_{\pi\in\Pi_\alpha}\theta_\pi$
(Thesis \S6.3):
(STAB) each $\theta_\alpha$ is stable;
(TRANS) $\theta_\alpha(y)=e(G_\alpha)\theta_1(x)$ for stably conjugate regular $x\in G(F)$, $y\in G_\alpha(F)$;
(WHITT) for every Whittaker datum of $G$ exactly one member of $\Pi_1$ has a Whittaker model.
\end{definition}

\section{The main theorem}
\begin{definition}[The multiplicity formula (statement)]
\label{MultiplicityFormula}
\lean{GSpinGGP.MultiplicityFormula}
\leanok
\uses{G, GGPData, GSpin, IsTempered, UnitaryRep, charG, mGeom, mult}
The statement of Thesis Thm.~7.2.2 (= Thm.~1.0.2): for every GGP triple and every irreducible tempered
$\pi$ with central character trivial on $A_{G,H}(F)$, $m(\pi)<\infty$ and $m(\pi)=\mgeom(\theta_\pi)$.
\end{definition}

\begin{theorem}[The multiplicity formula]
\label{multiplicity_formula}
\lean{GSpinGGP.multiplicity_formula}
\leanok
\uses{MultiplicityFormula}
The multiplicity formula holds (Thesis Thm.~7.2.2 = Thm.~1.0.2).
\end{theorem}
\begin{proof}
\leanok
\uses{multiplicityFormula_of_traceFormula}
By the strong induction on $\dim G$ of Thesis \S7.2, from the geometric expansion, the elliptic form of
Thesis Prop.~7.4.1(ii) and elliptic isolation (Chapter~\ref{chap:multformula}).
\end{proof}

\begin{theorem}[The count (Thesis (6.7.1))]
\label{sum_mult_eq_one}
\lean{GSpinGGP.sum_mult_eq_one}
\leanok
\uses{GGPData, InnerFormSystem, InnerTwist, IsQuasiSplit, TemperedPacketFamily, mult}
Let $G$ be quasi-split and $(\Pi_\alpha)$ as above, with finitely many inner forms. Then
$\sum_\alpha\sum_{\pi\in\Pi_\alpha} m(\pi)=1$.
\end{theorem}
\begin{proof}
\leanok
\uses{G, GSpin, IsTempered, LIT-germ_unique, LIT-hc_character_isQuasiCharacter, LIT-rodier_germ_one, UnitaryRep, WhittakerDatum, charG, germ, gspin, jHat, mGeom, multiplicity_formula, orbitChoice, orbitChoice_one_mem_nilReg, sum_mult_eq_germ_one}
By (6.7.7) the sum is the germ $c_{\theta_1,\mathcal O_1}(1)$. By Rodier's theorem applied to
the regular orbit $\mathcal O_1$, $c_{\theta_1,\mathcal O_1}(1)$ is the number of members of $\Pi_1$ with a Whittaker
model for the corresponding datum, which is $1$ by (WHITT). The multiplicity formula gives $m(\pi)<\infty$.
\end{proof}

\begin{theorem}[Main theorem: the local GGP conjecture for $\GSpin$]
\label{local_ggp}
\lean{GSpinGGP.local_ggp}
\leanok
\uses{GGPData, InnerFormSystem, InnerTwist, IsQuasiSplit, TemperedPacketFamily, mult}
(Thesis Thm.~6.4.1 = Thm.~1.0.1, $p$-adic case.) Let $G$ be quasi-split and $(\Pi_\alpha)_\alpha$ tempered
packets on the inner forms $G_\alpha$, $\alpha\in H^1(F,\SO(W))$, with central characters trivial on
$Z_{G,H}(F)$, satisfying (STAB), (TRANS), (WHITT). Then exactly one $\pi\in\bigsqcup_\alpha\Pi_\alpha$ has
$m(\pi)=1$, and $m(\pi')=0$ for all others.
\end{theorem}
\begin{proof}
\leanok
\uses{padic_quadratic_classes, sum_mult_eq_one}
The inner forms are finite in number (classification of $p$-adic quadratic forms). The sum of the
multiplicities is $1$, and each is a natural number, so exactly one is $1$ and the rest vanish
(Thesis \S6.7). Multiplicity one is not needed for this deduction.
\end{proof}


% Generated from the Lean dependency data; see the comment at the top of content.tex.
\chapter{Proof of the main theorem from the multiplicity formula}
\label{chap:proofmain}
Thesis \S6.5--6.7: summing the multiplicity formula over all inner forms and rearranging the geometric
side over stable classes in $\Gamma(G,H)$, everything cancels except the contribution of $x=1$.
For $x\in\GSpin(W)(F)$ write $W'_x=\ker(P(x)-1)$, $W''_x$ its orthogonal complement, and
$T_x\subset\GSpin(W)$ the centralizer of $x$ acting trivially on $W'_x$; $S_x=T_x/A_{G,H}$.

\section{\texorpdfstring{The space $\Gamma(G,H)$ and the geometric multiplicity}{The space Gamma(G,H) and the geometric multiplicity}}
\begin{definition}[Elements of $\Gamma(G,H)$]
\label{IsGammaElt}
\lean{GSpinGGP.GGPData.IsGammaElt}
\leanok
\uses{GGPData, GSpin}
$x\in\GSpin(W)(F)$ defines an element of $\Gamma(G,H)$ if $x$ is semisimple, $T_x$ is a torus and
$T_x(F)/A_{G,H}(F)$ is compact. This differs from Thesis \S6.5, which requires $T_x$ itself to be anisotropic.
\end{definition}

\begin{definition}[Classes of tori in $\Gamma(G,H)$]
\label{GammaClass}
\lean{GSpinGGP.GammaClass}
\leanok
\uses{GGPData, GSpin, IsGammaElt}
The index set of $\Gamma(G,H)$: $\GSpin(W)(F)$-conjugacy classes of the tori $T_x$, $x\in\Gamma(G,H)$.
\end{definition}

\begin{lemma}[Finiteness of the classes of tori]
\label{finite_gammaClass}
\lean{GSpinGGP.finite_gammaClass}
\leanok
\uses{GGPData, GammaClass}
$\Gamma(G,H)$ meets finitely many conjugacy classes of tori (Thesis \S6.5).
\end{lemma}
\begin{proof}
Not yet formalized. $T_x$ is determined by $W''_x$ (finitely many up to $\SO(W)$, by Witt) and a maximal torus
of $\SO(W''_x)$, of which there are finitely many conjugacy classes over a $p$-adic field.
\end{proof}

\begin{definition}[The function $\Delta$]
\label{deltaGGP}
\lean{GSpinGGP.deltaGGP}
\leanok
\uses{GGPData, GSpin, IsGammaElt}
$\Delta(x)=|\det((P(x)-1)|_{W''_x})|_F$ (Thesis (6.5.9)).
\end{definition}

\begin{definition}[The orbit $\mathcal O_x$]
\label{orbitChoice}
\lean{GSpinGGP.orbitChoice}
\leanok
\uses{G, GGPData, GSpin, IsGammaElt, gspin}
The regular nilpotent orbit $\mathcal O_x\in\Nilreg(\mathfrak g_x)$ selected by Luo's recipe (\cite{Luo} \S7.1)
applied to $(G'_x,H'_x)$; a junk value if there is none (then germs vanish anyway).
This differs from the choice in Thesis \S5.4: it is Luo's recipe applied to $(G'_x,H'_x)$.
\end{definition}

\begin{definition}[The geometric multiplicity $\mgeom$]
\label{mGeom}
\lean{GSpinGGP.mGeom}
\leanok
\uses{G, GGPData, GSpin, GammaClass, IsGammaElt, deltaGGP, germ, gspin, orbitChoice, weylDisc}
For $s>0$,
\[ m_{\mathrm{geom},s}(\theta)=\sum_{[T]}|W(T)|^{-1}\int_{T(F)/A_{G,H}(F)}D^G(t)^{1/2}c_{\theta,\mathcal O_t}(t)\Delta(t)^{s-1/2}\,dt, \]
the sum over the classes of tori of $\Gamma(G,H)$, with the probability Haar measure and
$W(T)=N_{\GSpin(W)(F)}(T)/Z_{\GSpin(W)(F)}(T)$; and $\mgeom(\theta)=\lim_{s\to0^+}m_{\mathrm{geom},s}(\theta)$
(Thesis \S7.2, Prop.~12.1.1). This differs from Thesis (6.5.3), which uses $N_{\GSpin(W'')(F)}(T)/T(F)$ for $W(T)$.
\end{definition}

\begin{proposition}[Linearity of $\mgeom$ on characters]
\label{mGeom_finset_sum}
\lean{GSpinGGP.mGeom_finset_sum}
\leanok
\uses{G, GGPData, GSpin, IsTempered, UnitaryRep, charG, mGeom}
For tempered $\pi_1,\dots,\pi_k$ with central characters trivial on $A_{G,H}(F)$,
$\mgeom(\sum_i\theta_{\pi_i})=\sum_i\mgeom(\theta_{\pi_i})$ (Thesis Prop.~12.1.1, used in (6.7.2)).
\end{proposition}
\begin{proof}
\leanok
\uses{CentralCharOn, GammaClass, IsGammaElt, IsQuasiCharacter, LIT-germ_unique, LIT-schur_central, QC, charG_mem_QC, deltaGGP, germ, gspin, jHat, mGeom_add, mult, orbitChoice, tendsto_mGeomS, weylDisc}
Each $\theta_{\pi_i}$ lies in $\QC(G,\eta_i)$ for its central character $\eta_i$, and $\mgeom$ is linear on the
span of these spaces (Chapter~\ref{chap:mgeom}).
\end{proof}

\begin{proposition}[Thesis (6.7.2)]
\label{packet_sum_mult}
\lean{GSpinGGP.packet_sum_mult}
\leanok
\uses{G, GGPData, InnerFormSystem, InnerTwist, TemperedPacketFamily, charG, mGeom, mult}
$\sum_{\pi\in\Pi_\alpha}m(\pi)=\mgeom(\theta_\alpha)$ for every inner form $\alpha$.
\end{proposition}
\begin{proof}
\leanok
\uses{mGeom_finset_sum, multiplicity_formula}
The multiplicity formula for $(G_\alpha,H_\alpha,\xi_\alpha)$ and linearity of $\mgeom$.
\end{proof}

\section{Germ transfer}
\begin{literature}[Transfer preserves germs]
\label{LIT-bp_germ_transfer}
\lean{GSpinGGP.bp_germ_transfer}
\leanok
\uses{G, GGPData, GSpin, InnerTwist, IsQuasiCharacter, IsQuasiSplit, StablyConjG, germ, gspin}
(Thesis (6.1.2).) Let $G$ be quasi-split, $\theta$ a stable quasi-character on $G(F)$ and $\theta'$ a stable
quasi-character on $G_\alpha(F)$ with $\theta'(y)=\theta(x)$ for stably conjugate regular $x,y$. Then for
semisimple, strongly stably conjugate $x\in\GSpin(W)(F)$, $y\in\GSpin(W_\alpha)(F)$,
$c_{\theta',\mathcal O'}(y)=c_{\theta,\mathcal O}(x)$ for all regular $\mathcal O,\mathcal O'$.

\assumedfrom{\cite{BP}, (12.1.2), stated there for pure inner forms; the proof is unchanged for these inner twists}
\end{literature}

\begin{proposition}[Germ transfer for packets]
\label{packet_germ_transfer}
\lean{GSpinGGP.packet_germ_transfer}
\leanok
\uses{G, GGPData, GSpin, InnerFormSystem, InnerTwist, IsQuasiSplit, StablyConjG, TemperedPacketFamily, charG, germ, gspin, kottwitzSign}
If $x\in\GSpin(W)(F)$, $y\in\GSpin(W_\alpha)(F)$ are semisimple and strongly stably conjugate, then
$c_{\theta_\alpha,\mathcal O'}(y)=e(G_\alpha)\,c_{\theta_1,\mathcal O}(x)$ for all regular $\mathcal O,\mathcal O'$
(Thesis \S6.7 after (6.7.4)).
\end{proposition}
\begin{proof}
\leanok
\uses{IsQuasiCharacter, IsStableFn, LIT-bp_germ_transfer, LIT-germ_unique, LIT-hc_character_isQuasiCharacter, jHat}
Apply the germ transfer to $e(G_\alpha)\theta_1$ and $\theta_\alpha$, using (STAB) and (TRANS).
\end{proof}

\section{\texorpdfstring{Stable conjugacy in $\Gamma(G,H)$}{Stable conjugacy in Gamma(G,H)}}
Thesis Prop.~6.6.1.

\begin{lemma}[Thesis Prop.~6.6.1(i), step 1]
\label{exists_isometry_fixedSub_of_qd}
\lean{GSpinGGP.exists_isometry_fixedSub_of_qd}
\leanok
\uses{GGPData, GSpin, InnerFormSystem, InnerTwist, IsGammaElt, IsQuasiSplit}
Let $G$ be quasi-split and $y\in\Gamma(G_\alpha,H_\alpha)$ with $G_{\alpha,y}$ quasi-split. Then $W'_{\alpha,y}$
embeds isometrically in $W$.
\end{lemma}
\begin{proof}
\leanok
\uses{padic_quadratic_classes, witt_cancellation, witt_decomposition}
By the classification of $p$-adic quadratic forms and Witt cancellation: if $W''_{\alpha,y}$ is $0$ or a
hyperbolic plane then $W_\alpha\cong W$; $\dim W''_{\alpha,y}=1$ is impossible; otherwise a second class of the
dimension and discriminant of $W''_{\alpha,y}$ exists.
\end{proof}

\begin{lemma}[Thesis Prop.~6.6.1(i), steps 2--3]
\label{exists_stronglyStablyConj_of_isometry}
\lean{GSpinGGP.exists_stronglyStablyConj_of_isometry}
\leanok
\uses{G, GGPData, GSpin, InnerFormSystem, InnerTwist, IsGammaElt, IsQuasiSplit, StablyConjG}
Given an isometric embedding $W'_{\alpha,y}\hookrightarrow W$, some $x\in\Gamma(G,H)$ is strongly stably
conjugate to $y$.
\end{lemma}
\begin{proof}
Not yet formalized. Transport the anisotropic torus through a cocycle in $S_y$ (\cite{BP} (12.1.6)), lift
$(1_{W'},\bar x'')$ from $\SO$ to $\GSpin$ and rescale by a scalar.
\end{proof}

\begin{proposition}[Thesis Prop.~6.6.1(i)]
\label{prop_12_5_1_i}
\lean{GSpinGGP.prop_12_5_1_i}
\leanok
\uses{G, GGPData, GSpin, InnerFormSystem, InnerTwist, IsGammaElt, IsQuasiSplit, StablyConjG}
Let $G$ be quasi-split and $y\in\Gamma(G_\alpha,H_\alpha)$ with $G_{\alpha,y}$ quasi-split. Then some
$x\in\Gamma(G,H)$ is strongly stably conjugate to $y$.
\end{proposition}
\begin{proof}
\leanok
\uses{exists_isometry_fixedSub_of_qd, exists_stronglyStablyConj_of_isometry}
Combine the two preceding lemmas.
\end{proof}

\begin{definition}[Torus transfer]
\label{TorusTransfer}
\lean{GSpinGGP.TorusTransfer}
\leanok
\uses{G, GGPData, GSpin, InnerFormSystem, InnerTwist, IsGammaElt, StablyConjG, deltaGGP, gspin, weylDisc}
For $x\in\Gamma(G,H)$, $y\in\Gamma(G_\alpha,H_\alpha)$: an $F$-rational isomorphism of topological groups
$T_x(F)\simeq T_y(F)$, the identity on $A_{G,H}$, with $x\mapsto y$, matching the regular parts $T_\natural$ and
the strong stable classes, and preserving $D^G$ and $\Delta$ (the conclusion of Thesis Prop.~6.6.1(ii)).
It need not preserve $|W(T)|$.
\end{definition}

\begin{lemma}[The torus isomorphism]
\label{exists_torusTransfer_core}
\lean{GSpinGGP.exists_torusTransfer_core}
\leanok
\uses{G, GGPData, GSpin, InnerFormSystem, InnerTwist, IsGammaElt, StablyConjG}
If $x\in\Gamma(G,H)$ and $y\in\Gamma(G_\alpha,H_\alpha)$ are strongly stably conjugate, then $\Ad(g)\circ\psi_\alpha$
restricts to an $F$-rational isomorphism of topological groups $T_x(F)\simeq T_y(F)$, the identity on scalars,
$x\mapsto y$, carrying $T_{x,\natural}$ onto $T_{y,\natural}$ with every $t$ stably conjugate to its image.
\end{lemma}
\begin{proof}
Not yet formalized. The values of $t\mapsto gC(\varphi_\alpha)(t)g^{-1}$ are Galois-fixed, hence rational
(Galois descent), commute with $y$ and fix $W'_y$.
\end{proof}

\begin{lemma}[$\Delta$ is a stable invariant]
\label{deltaGGP_eq_of_stablyConjVia}
\lean{GSpinGGP.deltaGGP_eq_of_stablyConjVia}
\leanok
\uses{GGPData, GSpin, InnerTwist, deltaGGP}
If $P(t')=hP(t)h^{-1}$ over $\Fbar$ with $h=P(g)\circ\varphi$, then $\Delta(t')=\Delta(t)$.
\end{lemma}
\begin{proof}
\leanok
\uses{IsGammaElt}
$h$ conjugates $(P(t)-1)|_{W''_t}$ to $(P(t')-1)|_{W''_{t'}}$; both determinants lie in $F$.
\end{proof}

\begin{lemma}[$D^G$ is a stable invariant]
\label{weylDisc_eq_of_stablyConjG}
\lean{GSpinGGP.weylDisc_eq_of_stablyConjG}
\leanok
\uses{G, GGPData, InnerTwist, StablyConjG, gspin, weylDisc}
If $t\in G(F)$ and $t'\in G_\alpha(F)$ are stably conjugate then $D^{G_\alpha}(t')=D^G(t)$.
\end{lemma}
\begin{proof}
\leanok
\uses{GSpin}
$\psi_\alpha\circ\Ad(g)$ conjugates $1-\Ad t$ to $1-\Ad t'$ over $\Fbar$; characteristic polynomials are unchanged by base change.
\end{proof}

\begin{proposition}[Thesis Prop.~6.6.1(ii)]
\label{prop_12_5_1_ii}
\lean{GSpinGGP.prop_12_5_1_ii}
\leanok
\uses{G, GGPData, GSpin, InnerFormSystem, InnerTwist, IsGammaElt, StablyConjG, TorusTransfer}
If $x\in\Gamma(G,H)$ and $y\in\Gamma(G_\alpha,H_\alpha)$ are strongly stably conjugate, there is a torus transfer
from $x$ to $y$. So the $\Gamma$-integrals over $T_x(F)/A(F)$ and $T_y(F)/A(F)$ match locally.
\end{proposition}
\begin{proof}
\leanok
\uses{deltaGGP, deltaGGP_eq_of_stablyConjVia, exists_torusTransfer_core, gspin, weylDisc, weylDisc_eq_of_stablyConjG}
The torus isomorphism, with the invariance of $\Delta$ and $D^G$.
\end{proof}

\section{The signed count}
Thesis (6.7.6) and Prop.~6.6.1(iii).

\begin{definition}[The counts $n_\alpha(x)$]
\label{n_alpha}
\lean{GSpinGGP.n_alpha}
\leanok
\uses{G, GGPData, GSpin, InnerFormSystem, InnerTwist, IsGammaElt, StablyConjG}
$n_\alpha(x)$ is the number of $\GSpin(W_\alpha)(F)$-conjugacy classes of $y\in\Gamma(G_\alpha,H_\alpha)$ strongly
stably conjugate to $x$ (Thesis \S6.7).
\end{definition}

\begin{definition}[The signed count]
\label{signedCount}
\lean{GSpinGGP.signedCount}
\leanok
\uses{GGPData, GSpin, InnerFormSystem, kottwitzSign, n_alpha}
$\sum_\alpha e(G_\alpha)\,n_\alpha(x)$ (Thesis (6.7.6)).
\end{definition}

\begin{literature}[Hilbert's Theorem 90 for $\mathrm{GL}_m$]
\label{LIT-hilbert90_GL}
\lean{GSpinGGP.hilbert90_GL}
\leanok
Every continuous cocycle $\sigma\mapsto a_\sigma\in\mathrm{GL}_m(\Fbar)$ of $\mathrm{Gal}(\Fbar/F)$ is a coboundary:
$a_\sigma=M^{-1}\sigma(M)$.

\assumedfrom{\cite{SerreLF}, Ch.~X \S1 Prop.~3; \cite{SerreCG}, Ch.~III \S1.1}
\end{literature}

\begin{lemma}[The map $q_x:H^1(F,S_x)\to\{\alpha\}$]
\label{existsUnique_twistIdx}
\lean{GSpinGGP.existsUnique_twistIdx}
\leanok
\uses{GGPData, GSpin, InnerFormSystem, InnerTwist}
Every class $\beta\in H^1(F,S_x)$ determines a unique inner form $\alpha=q_x(\beta)$: the twist of $W$ by a
cocycle representing $\beta$ is isometric to exactly one $W_\alpha$.
\end{lemma}
\begin{proof}
\leanok
\uses{LIT-hilbert90_GL}
Hilbert 90 writes the cocycle as $M^{-1}\sigma(M)$; $v\mapsto Q(M^{-1}v)$ is an $F$-form of the discriminant of $W$. Cohomologous cocycles give $F$-isometric forms by Galois descent.
\end{proof}

\begin{lemma}[$S_x$ is commutative]
\label{torusBarFbar_commute}
\lean{GSpinGGP.torusBarFbar_commute}
\leanok
\uses{GGPData, GSpin, InnerTwist, IsGammaElt}
For $x\in\Gamma(G,H)$, the group $S_x(\Fbar)$ is commutative.
\end{lemma}
\begin{proof}
Not yet formalized. $S_x$ is connected reductive with commutative $F$-points, which are Zariski dense.
\end{proof}

\begin{literature}[Finiteness of $H^1(F,S_x)$]
\label{LIT-finite_H1_torusBar}
\lean{GSpinGGP.finite_H1_torusBar}
\leanok
\uses{GGPData, GSpin, InnerTwist, IsGammaElt}
For $x\in\Gamma(G,H)$, $H^1(F,S_x)$ is finite.

\assumedfrom{\cite{SerreCG}, III \S4.3 Thm.~4 (Borel--Serre); Tate--Nakayama}
\end{literature}

\begin{lemma}[Thesis Prop.~6.6.1(iii)]
\label{n_alpha_eq_card_twistIdx}
\lean{GSpinGGP.n_alpha_eq_card_twistIdx}
\leanok
\uses{GGPData, GSpin, InnerFormSystem, InnerTwist, IsGammaElt, n_alpha}
For $x\in\Gamma(G,H)$ and every $\alpha$, $n_\alpha(x)=|q_x^{-1}(\alpha)|$, $q_x^{-1}(\alpha)\subset H^1(F,S_x)$.
This differs from Thesis Prop.~6.6.1(iii), which uses $H^1(F,T_x)$.
\end{lemma}
\begin{proof}
Not yet formalized. Galois-cohomological parametrization of stable classes, then Witt cancellation.
\end{proof}

\begin{literature}[The Kottwitz sign is multiplicative on $H^1(F,S_x)$]
\label{LIT-kottwitzSign_twistIdx_mul}
\lean{GSpinGGP.kottwitzSign_twistIdx_mul}
\leanok
\uses{GGPData, GSpin, InnerFormSystem, InnerTwist, IsGammaElt, IsQuasiSplit, kottwitzSign}
For $G$ quasi-split and $x\in\Gamma(G,H)$, $\beta\mapsto e(G_{q_x(\beta)})$ is multiplicative on $H^1(F,S_x)$.

\assumedfrom{\cite{Kot83}; \cite{BP}, \S12.1, (12.1.6), via Tate--Nakayama}
\end{literature}

\begin{lemma}[The sign takes the value $-1$]
\label{exists_twistIdx_kottwitzSign_neg}
\lean{GSpinGGP.exists_twistIdx_kottwitzSign_neg}
\leanok
\uses{GGPData, GSpin, InnerFormSystem, InnerTwist, IsGammaElt, IsQuasiSplit, kottwitzSign}
For $G$ quasi-split and $x\in\Gamma(G,H)$ not central, some $\beta\in H^1(F,S_x)$ has $e(G_{q_x(\beta)})=-1$.
\end{lemma}
\begin{proof}
Not yet formalized. By \cite{BP} (12.1.6) it suffices that the inverse image of $S_x$ in the simply connected cover is connected; its character group is torsion free.
\end{proof}

\begin{theorem}[The signed count]
\label{signed_count}
\lean{GSpinGGP.signed_count}
\leanok
\uses{GGPData, GSpin, InnerFormSystem, IsGammaElt, IsQuasiSplit, signedCount}
Let $G$ be quasi-split and $x\in\Gamma(G,H)$ with $G_x$ quasi-split. Then
$\sum_\alpha e(G_\alpha)n_\alpha(x)=1$ if $x$ is central and $0$ otherwise.
This replaces the Kottwitz-sign cancellation of Thesis \S6.7 after (6.7.6).
\end{theorem}
\begin{proof}
\leanok
\uses{InnerTwist, LIT-finite_H1_torusBar, LIT-kottwitzSign_twistIdx_mul, existsUnique_twistIdx, exists_twistIdx_kottwitzSign_neg, gspin, kottwitzSign, n_alpha, n_alpha_eq_card_twistIdx, padic_quadratic_classes, torusBarFbar_commute}
The sum equals $\sum_{\beta\in H^1(F,S_x)}e(G_{q_x(\beta)})$. If $x$ is central this is $e(G)=1$. Otherwise the
sign is a character of the finite group $H^1(F,S_x)$ taking the value $-1$, so the sum vanishes.
\end{proof}

\section{Rearrangement and the atom}
\begin{lemma}[The stable rearrangement at $s>0$]
\label{stable_rearrangement_pos}
\lean{GSpinGGP.stable_rearrangement_pos}
\leanok
\uses{G, GGPData, GSpin, GammaClass, InnerFormSystem, InnerTwist, IsGammaElt, IsQuasiSplit, StablyConjG, TemperedPacketFamily, TorusTransfer, charG, deltaGGP, germ, gspin, kottwitzSign, n_alpha, orbitChoice, signedCount, weylDisc}
Given germ transfer, vanishing of the non-quasi-split classes, and Prop.~6.6.1(i),(ii): for $s>0$,
\[ \sum_\alpha m_{\mathrm{geom},s}(\theta_\alpha)=\int_{\Gamma(G,H)}\frac{\sum_\alpha e(G_\alpha)n_\alpha(x)}{n_1(x)}
c_{\theta_1}(x)D^G(x)^{1/2}\Delta(x)^{s-1/2}\,dx \] (Thesis (6.7.2)--(6.7.5) before the limit).
\end{lemma}
\begin{proof}
Not yet formalized. At fixed $s>0$ the integrals converge absolutely; $y\mapsto x$ is a measure-preserving
injection $\Gamma^{qd}_{\mathrm{stab}}(G_\alpha,H_\alpha)\hookrightarrow\Gamma^{qd}_{\mathrm{stab}}(G,H)$, and
(TRANS) gives $c_{\theta_\alpha}(y)=e(G_\alpha)c_{\theta_1}(x)$.
\end{proof}

\begin{proposition}[Thesis (6.7.5)]
\label{stable_rearrangement}
\lean{GSpinGGP.stable_rearrangement}
\leanok
\uses{G, GGPData, GSpin, GammaClass, InnerFormSystem, InnerTwist, IsGammaElt, IsQuasiSplit, StablyConjG, TemperedPacketFamily, TorusTransfer, charG, deltaGGP, germ, gspin, kottwitzSign, mGeom, n_alpha, orbitChoice, signedCount, weylDisc}
Under the same hypotheses, \[ \sum_\alpha\mgeom(\theta_\alpha)=\lim_{s\to0^+}\int_{\Gamma(G,H)}
\frac{\sum_\alpha e(G_\alpha)n_\alpha(x)}{n_1(x)}c_{\theta_1}(x)D^G(x)^{1/2}\Delta(x)^{s-1/2}\,dx. \]
\end{proposition}
\begin{proof}
\leanok
\uses{CentralCharOn, IsTempered, LIT-schur_central, QC, UnitaryRep, charG_mem_QC, mult, padic_quadratic_classes, stable_rearrangement_pos, tendsto_mGeomS}
The rearrangement at $s>0$, and the finite sum over $\alpha$ commutes with the limit.
\end{proof}

\begin{lemma}[Central translation of the germ]
\label{germ_baseChar_central}
\lean{GSpinGGP.germ_baseChar_central}
\leanok
\uses{G, GGPData, GSpin, InnerFormSystem, InnerTwist, TemperedPacketFamily, charG, germ, gspin, orbitChoice}
For central $c\in A_{G,H}(F)$, $c_{\theta_1,\mathcal O_c}(c)=c_{\theta_1,\mathcal O_1}(1)$.
\end{lemma}
\begin{proof}
\leanok
\uses{IsQuasiCharacter, IsTempered, LIT-germ_unique, LIT-hc_character_exists, LIT-hc_character_isQuasiCharacter, LIT-hc_regularSet_open_conull, UnitaryRep, jHat}
$\theta_1(cy)=\theta_1(y)$ on $\Greg(F)$ since the central characters are trivial on $Z_{G,H}(F)$; $e^X$ is
regular for small regular $X$, so the expansion at $1$ is one at $c$ (uniqueness of germs); $\mathcal O_c=\mathcal O_1$.
\end{proof}

\begin{lemma}[The atom]
\label{atom_value}
\lean{GSpinGGP.atom_value}
\leanok
\uses{G, GGPData, GSpin, InnerFormSystem, InnerTwist, IsGammaElt, IsQuasiSplit, TemperedPacketFamily, charG, deltaGGP, germ, gspin, n_alpha, orbitChoice, signedCount, weylDisc}
Assuming the signed count, the contribution of the class of $T=A_{G,H}$ to the stable integral is
$c_{\theta_1,\mathcal O_1}(1)$: $|W(T)|=1$, $T(F)/A(F)$ is a point, the weight is $1$ and $D^G=\Delta=1$.
\end{lemma}
\begin{proof}
\leanok
\uses{StablyConjG, germ_baseChar_central}
Direct computation, with central translation of the germ.
\end{proof}

\begin{proposition}[Only the atom contributes]
\label{mGeomStabS_eq_germ_one}
\lean{GSpinGGP.mGeomStabS_eq_germ_one}
\leanok
\uses{G, GGPData, GSpin, GammaClass, InnerFormSystem, InnerTwist, IsGammaElt, IsQuasiSplit, TemperedPacketFamily, charG, deltaGGP, germ, gspin, n_alpha, orbitChoice, signedCount, weylDisc}
If the signed count holds on quasi-split classes, then for $s>0$ the stable integral equals
$c_{\theta_1,\mathcal O_1}(1)$ (Thesis \S6.7 after (6.7.6)).
\end{proposition}
\begin{proof}
\leanok
\uses{LIT-nilReg_centralizer_nonempty_imp, atom_value, jHat}
Non-central and non-quasi-split classes have weight $0$; the atom gives the value above.
\end{proof}

\begin{lemma}[Luo's orbit at $x=1$]
\label{exists_orbitSpec_one_of_even}
\lean{GSpinGGP.exists_orbitSpec_one_of_even}
\leanok
\uses{G, GGPData, GSpin, IsQuasiSplit, gspin}
Let $G$ be quasi-split with $\dim V$ or $\dim W$ even and $\ge 4$. Then some $\mathcal O\in\Nilreg(\mathfrak g)$
satisfies Luo's condition at $x=1$ (\cite{Luo} \S7.1).
\end{lemma}
\begin{proof}
\leanok
\uses{LIT-exists_regularNil_inOrbitNu, LIT-nilReg_nonempty_iff_quasiSplit}
Luo's orbit in the even factor, paired with a regular nilpotent of the other factor.
\end{proof}

\begin{lemma}[$\mathcal O_1$ is regular]
\label{orbitChoice_one_mem_nilReg}
\lean{GSpinGGP.orbitChoice_one_mem_nilReg}
\leanok
\uses{G, GGPData, GSpin, IsQuasiSplit, gspin, orbitChoice}
For $G$ quasi-split, $\mathcal O_1\in\Nilreg(\mathfrak g)$.
\end{lemma}
\begin{proof}
\leanok
\uses{IsGammaElt, LIT-nilReg_centralizer_nonempty_iff, exists_orbitSpec_one_of_even}
Luo's recipe applied to $(G,H)$ selects a regular orbit, which exists since $G$ is quasi-split.
\end{proof}

\begin{theorem}[Thesis (6.7.7)]
\label{sum_mult_eq_germ_one}
\lean{GSpinGGP.sum_mult_eq_germ_one}
\leanok
\uses{G, GGPData, GSpin, InnerFormSystem, InnerTwist, IsQuasiSplit, TemperedPacketFamily, charG, germ, gspin, mult, orbitChoice}
Let $G$ be quasi-split. Then $\sum_\alpha\sum_{\pi\in\Pi_\alpha}m(\pi)=c_{\theta_1,\mathcal O_1}(1)$.
\end{theorem}
\begin{proof}
\leanok
\uses{GammaClass, IsGammaElt, LIT-nilReg_centralizer_nonempty_imp, StablyConjG, TorusTransfer, deltaGGP, jHat, kottwitzSign, mGeom, mGeomStabS_eq_germ_one, n_alpha, packet_germ_transfer, packet_sum_mult, prop_12_5_1_i, prop_12_5_1_ii, signedCount, signed_count, stable_rearrangement, weylDisc}
Sum (6.7.2) over $\alpha$; rearrange by (6.7.5) using germ transfer and Prop.~6.6.1(i),(ii); the
signed count kills every class but the atom.
\end{proof}


% Generated from the Lean dependency data; see the comment at the top of content.tex.
\chapter{The multiplicity formula}
\label{chap:multformula}
Thesis Ch.~7. The multiplicity formula is proved by strong induction on $\dim G$, together with the
geometric expansion of the distribution $J$. Test functions are taken in $C_c^\infty$ modulo $A_G(F)$
instead of the Harish-Chandra Schwartz space; every stated theorem of the trace formula uses this space.

\section{\texorpdfstring{Strongly cuspidal functions and the distribution $J$}{Strongly cuspidal functions and the distribution J}}
\begin{definition}[Strongly cuspidal test functions]
\label{ScuspTestFn}
\lean{GSpinGGP.ScuspTestFn}
\leanok
\uses{CentralCharOn, G, GGPData, GSpin, H, N}
$\scusp(A_G(F)\backslash G(F),\chi^{-1})$, $C_c^\infty$ version: locally constant $f$ on $G(F)$ with
$f(ag)=\chi(a)^{-1}f(g)$ for $a\in A_G(F)$, compactly supported modulo $A_G(F)$, and strongly cuspidal
($\int_{U(F)}f(pu)\,du=0$ for every proper parabolic $P=MU$ and $p\in P(F)$) (Thesis Ch.~1, \S4.5, \S7.1).
\end{definition}

\begin{definition}[The distribution $J$]
\label{distribJ}
\lean{GSpinGGP.distribJ}
\leanok
\uses{G, GGPData, GSpin, H, N, xi}
\[ J(f)=\int_{H(F)\backslash G(F)/A_G(F)}K(f,x)\,dx,\qquad K(f,x)=\int_{H(F)/A_{G,H}(F)}f(x^{-1}hx)\xi(h)\,dh \]
(Thesis \S7.1, \S10.1), for the Weil quotient measure on $H(F)\backslash G(F)/A_G(F)$.
\end{definition}

\begin{theorem}[Convergence of $J$]
\label{distribJ_convergent}
\lean{GSpinGGP.GGPData.distribJ_convergent}
\leanok
\uses{CentralCharOn, G, GGPData, GSpin, H, N, ScuspTestFn, xi}
For $f\in\scusp(A_G\backslash G,\chi^{-1})$, $x\mapsto K(f,x)$ is integrable on $H(F)\backslash G(F)/A_G(F)$
(Thesis Thm.~10.1.1(ii) with Prop.~8.4.1(iii); = Thm.~1.0.3).
\end{theorem}
\begin{proof}
Not yet formalized. The bound $|K(f,x)|\le\nu(f)\,\Xi^{H\backslash G}(x)^2\sigma(x)^{-d}$ for strongly cuspidal $f$,
and integrability of $\Xi^{H\backslash G}(x)^2\sigma(x)^{-d}$ for $d$ large.
\end{proof}

\begin{lemma}[$J$ is linear]
\label{distribJ_add}
\lean{GSpinGGP.GGPData.distribJ_add}
\leanok
\uses{CentralCharOn, G, GGPData, ScuspTestFn, distribJ}
$J(f+g)=J(f)+J(g)$ for $f,g\in\scusp(A_G\backslash G,\chi^{-1})$ (Thesis \S10.1).
\end{lemma}
\begin{proof}
\leanok
\uses{GSpin, H, N, distribJ_convergent, xi}
Linearity of the integral, given convergence.
\end{proof}

\begin{definition}[The quasi-character $\theta_f$]
\label{thetaF}
\lean{GSpinGGP.thetaF}
\leanok
\uses{G, GGPData, GSpin, H, N, gspin}
For strongly cuspidal $f$, $\theta_f$ is the function representing the orbital pairing
$f'\mapsto\int_{G(F)/A_G(F)}\int_{G(F)}f(x^{-1}gx)f'(g)\,dg\,dx$ (Thesis \S4.8, \cite{BP} \S5.5);
$0$ if there is none. Convention: $\theta_f(ag)=\chi(a)^{-1}\theta_f(g)$.
\end{definition}

\begin{literature}[$\theta_f$ represents the orbital pairing]
\label{LIT-bp_thetaF_exists}
\lean{GSpinGGP.LitTF.bp_thetaF_exists}
\leanok
\uses{G, GGPData, GSpin, IsQuasiCharacter, gspin}
Let $f_0\in C_c^\infty(G(F))$ be strongly cuspidal. Then there is a quasi-character $\theta$, locally integrable
and locally constant on $\Greg(F)$, with $\int_{G(F)}\theta f'\,dg=\int_{G(F)/A_G(F)}\int_{G(F)}f_0(x^{-1}gx)f'(g)\,dg\,dx$
for every $f'\in C_c^\infty(G(F))$, the outer integrand being integrable.

\assumedfrom{\cite{BP}, Thm.~5.5.1(ii),(iv), Prop.~5.6.1(ii) (= Thesis Thm.~4.7.1, Prop.~4.8.1(ii)); \cite{Art91}; Weyl integration formula}
\end{literature}

\begin{lemma}[$\chi$-transfer of $\theta_f$]
\label{isThetaF_chiProj}
\lean{GSpinGGP.isThetaF_chiProj}
\leanok
\uses{CentralCharOn, G, GGPData, GSpin, H, IsQuasiCharacter, N, gspin}
For strongly cuspidal $f_0\in C_c^\infty(G(F))$, the orbital pairing of its $\chi$-projection $P_\chi f_0$
is represented by a quasi-character (Thesis \S4.8, \S7.1).
\end{lemma}
\begin{proof}
\leanok
\uses{LIT-bp_thetaF_exists, ScuspTestFn, jHat}
Glue central translates of the quasi-character of a cut-off of $P_\chi f_0$ along the fibres of the Harish-Chandra map $H_G$.
\end{proof}

\begin{proposition}[$\theta_f$ of a strongly cuspidal function]
\label{thetaF_spec}
\lean{GSpinGGP.thetaF_spec}
\leanok
\uses{CentralCharOn, G, GGPData, GSpin, H, IsQuasiCharacter, N, ScuspTestFn, gspin, thetaF}
For $f\in\scusp(A_G\backslash G,\chi^{-1})$, $\theta_f$ represents the orbital pairing of $f$ and is a
quasi-character (Thesis \S4.8, Prop.~4.8.1(ii)).
\end{proposition}
\begin{proof}
\leanok
\uses{LIT-hc_regularSet_open_conull, isThetaF_chiProj, jHat, weylDisc}
From the $\chi$-transfer and uniqueness on the open conull regular set.
\end{proof}

\begin{definition}[$\QC(G(F),\chi^{-1})$]
\label{QC}
\lean{GSpinGGP.QC}
\leanok
\uses{CentralCharOn, G, GGPData, GSpin, IsQuasiCharacter, gspin, jHat}
Quasi-characters $\theta$ on $G(F)$ with $\theta(ag)=\chi(a)^{-1}\theta(g)$ for $a\in A_G(F)$, $g\in\Greg(F)$
(Thesis \S4.1.1, \S7.2).
\end{definition}

\begin{lemma}[$\theta_f\in\QC(G,\chi^{-1})$]
\label{thetaF_mem_QC}
\lean{GSpinGGP.thetaF_mem_QC}
\leanok
\uses{CentralCharOn, G, GGPData, QC, ScuspTestFn, thetaF}
For $f\in\scusp(A_G\backslash G,\chi^{-1})$, $\theta_f\in\QC(G(F),\chi^{-1})$ (Thesis Prop.~4.8.1(ii)).
\end{lemma}
\begin{proof}
\leanok
\uses{GSpin, H, IsQuasiCharacter, LIT-hc_regularSet_open_conull, N, gspin, thetaF_spec}
From the previous proposition and central equivariance of $f$.
\end{proof}

\begin{lemma}[Characters are quasi-characters]
\label{charG_mem_QC}
\lean{GSpinGGP.charG_mem_QC}
\leanok
\uses{CentralCharOn, G, GGPData, GSpin, IsTempered, QC, UnitaryRep, charG, mult, orbitChoice}
If $\pi$ is irreducible tempered with central character $\eta$ on $A_G(F)$ then $\theta_\pi\in\QC(G(F),\eta)$
(Thesis Ch.~1, \S4.4).
\end{lemma}
\begin{proof}
\leanok
\uses{IsQuasiCharacter, LIT-hc_character_exists, LIT-hc_character_isQuasiCharacter, LIT-hc_regularSet_open_conull, gspin}
Harish-Chandra's local expansion, and central translation of characters.
\end{proof}

\section{The induction}
Thesis \S7.2 and \S7.4.

\begin{definition}[The geometric expansion at $(G,\chi)$]
\label{GeomAt}
\lean{GSpinGGP.GeomAt}
\leanok
\uses{CentralCharOn, G, GGPData, ScuspTestFn, distribJ, mGeom, thetaF}
$J(f)=\mgeom(\theta_f)$ for every $f\in\scusp(A_G\backslash G,\chi^{-1})$ (Thesis Thm.~7.2.1 = Thm.~12.2.1 = Thm.~1.0.5).
\end{definition}

\begin{definition}[The multiplicity formula at one triple]
\label{MultAt}
\lean{GSpinGGP.MultAt}
\leanok
\uses{G, GGPData, GSpin, IsTempered, UnitaryRep, charG, mGeom, mult}
For every irreducible tempered $\pi$ of $G(F)$ trivial on $A_{G,H}(F)$: $m(\pi)<\infty$ and
$m(\pi)=\mgeom(\theta_\pi)$.
\end{definition}

\begin{definition}[The induction hypothesis (HYP)]
\label{HYP}
\lean{GSpinGGP.HYP}
\leanok
\uses{CentralCharOn, GGPData, GeomAt, MultAt}
The geometric expansion (for every central character) and the multiplicity formula hold for every GGP
triple $(G',H',\xi')$ with $\dim G'<\dim G$ (Thesis \S7.2).
\end{definition}

\begin{definition}[Virtual tempered representations]
\label{VirtualTemp}
\lean{GSpinGGP.VirtualTemp}
\leanok
\uses{CentralCharOn, G, GGPData, GSpin, IsTempered, UnitaryRep, charG, mult, orbitChoice}
$\Rtemp(G,\eta)\otimes\mathbb C$: finite formal linear combinations of classes of irreducible tempered
representations with central character $\eta$ on $A_G(F)$ (Thesis \S3.4, \S7.4), with the linear extensions
$\theta_\Pi$, $m(\Pi)$, $\mgeom(\Pi)$.
\end{definition}

\begin{definition}[The defect]
\label{defect}
\lean{GSpinGGP.defect}
\leanok
\uses{CentralCharOn, G, GGPData, GSpin, IsTempered, UnitaryRep, VirtualTemp, charG, mGeom, mult, orbitChoice}
$\Pi\mapsto m(\Pi)-\mgeom(\Pi)$, extended linearly (Thesis Prop.~7.4.1(ii)). The multiplicity formula
for $\pi$ is $\mathrm{defect}[\pi]=0$.
\end{definition}

\begin{literature}[Elliptic isolation]
\label{LIT-bp_scusp_surj_ell}
\lean{GSpinGGP.LitTF.bp_scusp_surj_ell}
\leanok
\uses{G, GGPData, GSpin, IsQuasiCharacter, gspin}
There is an open, closed, invariant, $A_G(F)$-stable $\Omega\subset G(F)$ containing the elliptic regular set
and compact modulo conjugation and $A_G(F)$, such that every quasi-character supported (compactly
modulo conjugation) in $\Omega$ is $\theta_{f_0}$ on $\Greg(F)$ for some strongly cuspidal $f_0\in C_c^\infty(G(F))$
supported in $\Omega$.

\assumedfrom{\cite{BP}, Cor.~5.7.2(i) (= Thesis Cor.~4.8.3(i)), stated there for compact centre; adapted to a nontrivial split centre $A_G$ by its local proof}
\end{literature}

\begin{proposition}[Elliptic isolation, $\chi$-equivariant]
\label{exists_scusp_ell}
\lean{GSpinGGP.exists_scusp_ell}
\leanok
\uses{CentralCharOn, G, GGPData, GSpin, IsTempered, ScuspTestFn, UnitaryRep, VirtualTemp, charG, gspin, mult, orbitChoice, thetaF}
For every virtual tempered $\Pi$ with central character $\eta$ there is $f\in\scusp(A_G\backslash G,\eta)$ with
$\theta_f=\theta_\Pi$ on $\ellreg(F)$ (Thesis Cor.~4.8.3(i) in $\chi$-equivariant form).
This differs from the Thesis, which uses \cite{BP} Cor.~5.7.2(iv) and Schwartz functions.
\end{proposition}
\begin{proof}
\leanok
\uses{H, IsQuasiCharacter, LIT-bp_scusp_surj_ell, LIT-bp_thetaF_exists, LIT-hc_regularSet_open_conull, N, QC, charG_mem_QC, jHat}
Cut $\theta_\Pi$ off near the elliptic set, apply elliptic isolation, and project with $P_\chi$.
\end{proof}

\begin{proposition}[Thesis Prop.~7.4.1(ii), elliptic form]
\label{prop_11_5_1_ii_ell}
\lean{GSpinGGP.prop_11_5_1_ii_ell}
\leanok
\uses{CentralCharOn, G, GGPData, GSpin, HYP, IsTempered, ScuspTestFn, UnitaryRep, VirtualTemp, charG, defect, distribJ, gspin, mGeom, mult, orbitChoice, thetaF}
Assume (HYP). Let $f\in\scusp(A_G\backslash G,\eta)$ and $\Pi$ virtual tempered with central character $\eta$ such
that $\theta_\Pi=\theta_f$ on $\ellreg(F)$. Then $J(f)-\mgeom(\theta_f)=m(\Pi)-\mgeom(\Pi)$.
\end{proposition}
\begin{proof}
\leanok
\uses{defect_eq_zero_of_vChar_eq_zero_ell, exists_virtual_spectral_ell}
The spectral expansion gives a $\Pi_f$ with the same elliptic character and
$J(f)-\mgeom(\theta_f)=\mathrm{defect}(\Pi_f)$; $\Pi-\Pi_f$ has vanishing elliptic character, so it lies in $\Rind$,
where the defect vanishes (Chapter~\ref{chap:spectral}).
\end{proof}

\begin{proposition}[Geometric expansion implies multiplicity formula]
\label{multAt_of_hyp_geom}
\lean{GSpinGGP.multAt_of_hyp_geom}
\leanok
\uses{CentralCharOn, GGPData, GeomAt, HYP, MultAt}
(Thesis Prop.~7.4.1(iii), direction Thm.~7.2.1 $\Rightarrow$ Thm.~7.2.2.) Under (HYP), if the geometric
expansion holds at $G$ for all central characters, the multiplicity formula holds at $G$.
\end{proposition}
\begin{proof}
\leanok
\uses{G, GSpin, H, IsTempered, LIT-emoryTakeda_mult_le_one, LIT-schur_central, ScuspTestFn, UnitaryRep, VirtualTemp, charG, defect, distribJ, exists_scusp_ell, gspin, mGeom, mult, orbitChoice, prop_11_5_1_ii_ell, thetaF, xi}
Let $\eta$ be the central character of $\pi$ (Schur) and choose $f$ with $\theta_f=\theta_\pi$ on the elliptic set
(elliptic isolation). Then $0=J(f)-\mgeom(\theta_f)=\mathrm{defect}[\pi]$. Finiteness is multiplicity one.
\end{proof}

\begin{theorem}[The geometric expansion under (HYP)]
\label{geomAt_of_hyp}
\lean{GSpinGGP.geomAt_of_hyp}
\leanok
\uses{CentralCharOn, GGPData, GeomAt, HYP}
(Thesis Thm.~12.2.1 under (HYP).) If (HYP) holds for $G$, then $J(f)=\mgeom(\theta_f)$ for every central
character $\chi$ and every $f\in\scusp(A_G\backslash G,\chi^{-1})$.
\end{theorem}
\begin{proof}
\leanok
\uses{G, GSpin, GammaClass, H, IsGammaElt, IsQuasiCharacter, LIT-germ_unique, LIT-hc_regularSet_open_conull, N, QC, ScuspTestFn, deltaGGP, distribJ, distribJ_add, geomAt_away, geomAt_near, germ, gspin, isThetaF_chiProj, jHat, mGeom, mGeom_add, mem_mul_parabolicU_iff, orbitChoice, thetaF, thetaF_mem_QC, thetaF_spec, weylDisc}
Choose the invariant clopen $\Omega\ni1$ of the expansion near $1$; cut $f=f_1+(f-f_1)$ with $f_1$ supported in $\Omega$
and $\theta_{f_1}=1_\Omega\theta_f$; apply the expansion near $1$ to $f_1$ and the expansion away from $1$ to
$f-f_1$, and add, using linearity of $J$ and $\mgeom$ (Chapter~\ref{chap:geometric}).
\end{proof}

\begin{theorem}[The induction]
\label{geom_mult_all}
\lean{GSpinGGP.geom_mult_all}
\leanok
\uses{CentralCharOn, GGPData, GeomAt, MultAt}
For every GGP triple, the geometric expansion (for every central character) and the multiplicity formula hold.
\end{theorem}
\begin{proof}
\leanok
\uses{HYP, geomAt_of_hyp, multAt_of_hyp_geom}
Strong induction on $\dim G$: (HYP) is the induction hypothesis; apply the two preceding results.
\end{proof}

\begin{theorem}[Multiplicity formula from the trace formula]
\label{multiplicityFormula_of_traceFormula}
\lean{GSpinGGP.multiplicityFormula_of_traceFormula}
\leanok
\uses{MultiplicityFormula}
$m(\pi)=\mgeom(\theta_\pi)$ (and $m(\pi)<\infty$) for every GGP triple and every tempered $\pi$ trivial on
$A_{G,H}(F)$ (Thesis Thm.~7.2.2 = Thm.~1.0.2).
\end{theorem}
\begin{proof}
\leanok
\uses{CentralCharOn, GGPData, GeomAt, MultAt, geom_mult_all}
The multiplicity-formula half of the induction.
\end{proof}


% Generated from the Lean dependency data; see the comment at the top of content.tex.
\chapter{The spectral side}
\label{chap:spectral}
Thesis Ch.~9--10 and Prop.~7.4.1: the spectral expansion of $J$, in the elliptic form used by the induction,
and the vanishing of the defect on properly induced representations.

\section{Parabolic induction}
\begin{definition}[Normalized parabolic induction]
\label{parInd}
\lean{GSpinGGP.GGPData.parInd}
\leanok
\uses{G, GGPData, GSpin, LIT-haar_G_isMulRightInvariant, LIT-integral_eq_mul_integral_compact_integral, UnitaryRep}
$i_P^G(\sigma)$: the completion of the smooth $f:G(F)\to H_\sigma$ with $f(pg)=\delta_P(p)^{1/2}\sigma(p)f(g)$ for
$\|f\|^2=\int_K\|f(k)\|^2dk$, with $G(F)$ acting by right translation (Thesis \S3.2).
\end{definition}

\begin{definition}[Properly induced virtual representations]
\label{Rind}
\lean{GSpinGGP.GGPData.Rind}
\leanok
\uses{CentralCharOn, G, GGPData, GSpin, IsTempered, UnitaryRep, VirtualTemp, charG, gspin, mult, orbitChoice, parInd}
$\Rind(G,\eta)$: the span of the virtual tempered representations whose character on $\Greg(F)$ is that of some
$i_P^G(\sigma)$, $P$ proper, $\sigma$ tempered with central character $\eta$ (Thesis \S7.4).
\end{definition}

\begin{literature}[Van Dijk's formula]
\label{LIT-vanDijk_charG_parInd_ell}
\lean{GSpinGGP.vanDijk_charG_parInd_ell}
\leanok
\uses{CentralCharOn, G, GGPData, GSpin, IsTempered, UnitaryRep, charG, gspin, parInd}
For $P$ proper and $\sigma$ tempered, $\theta_{i_P^G\sigma}$ vanishes on $\ellreg(F)$.

\assumedfrom{\cite{vD72}; \cite{BP}, (3.4.2)/(3.4.3) = Thesis (3.5.3.2)}
\end{literature}

\begin{literature}[Multiplicity of an induced representation]
\label{LIT-hc_vMult_eq_mult_parInd}
\lean{GSpinGGP.hc_vMult_eq_mult_parInd}
\leanok
\uses{CentralCharOn, G, GGPData, GSpin, IsTempered, UnitaryRep, VirtualTemp, charG, gspin, mult, orbitChoice, parInd}
If a virtual representation $\Pi$ has the character of $i_P^G(\sigma)$ on $\Greg(F)$, then $m(\Pi)=m(i_P^G\sigma)$.

\assumedfrom{Harish-Chandra (linear independence of characters); \cite{WalPl}, \S III.4, \S VII; \cite{BP}, \S2.2, \S2.7}
\end{literature}

\begin{literature}[Vanishing elliptic character means properly induced]
\label{LIT-arthur_mem_Rind_of_vChar_ell}
\lean{GSpinGGP.arthur_mem_Rind_of_vChar_ell}
\leanok
\uses{CentralCharOn, G, GGPData, GSpin, IsTempered, Rind, UnitaryRep, VirtualTemp, charG, gspin, mult, orbitChoice}
If $\Pi\in\Rtemp(G,\eta)$ has $\theta_\Pi=0$ on $\ellreg(F)$, then $\Pi\in\Rind(G,\eta)$.

\assumedfrom{\cite{Art93}, \S6 (orthogonality of elliptic tempered characters); \cite{BP}, \S2.7, (2.7.1)}
\end{literature}

\begin{literature}[$\theta_f$ is a virtual character on the elliptic set]
\label{LIT-bp_exists_vChar_eq_thetaF_ell}
\lean{GSpinGGP.bp_exists_vChar_eq_thetaF_ell}
\leanok
\uses{CentralCharOn, G, GGPData, GSpin, IsTempered, ScuspTestFn, UnitaryRep, VirtualTemp, charG, gspin, mult, orbitChoice, thetaF}
For $f\in\scusp(A_G\backslash G,\eta)$ there is a virtual tempered $\Pi$ with central character $\eta$ with
$\theta_\Pi=\theta_f$ on $\ellreg(F)$.

\assumedfrom{\cite{BP}, Prop.~5.6.1(ii), Lemma 5.4.2(i) (= Thesis Prop.~4.8.1), adapted to the $\chi$-equivariant setting and a finite sum}
\end{literature}

\section{The spectral expansion}
\begin{proposition}[Thesis Prop.~7.4.1(i) at an induced representation]
\label{mult_parInd_eq_mGeom}
\lean{GSpinGGP.mult_parInd_eq_mGeom}
\leanok
\uses{CentralCharOn, G, GGPData, GSpin, HYP, IsTempered, UnitaryRep, charG, mGeom, mult, parInd}
Assume (HYP). For $\pi=i_Q^G(\sigma)$, $Q$ proper and $\sigma$ tempered with central character $\eta$,
$m(\pi)=\mgeom(\theta_\pi)$.
\end{proposition}
\begin{proof}
Not yet formalized. $m(i_Q^G\sigma)=m(\tilde\sigma)$ for the GGP triple of the Levi (Thesis Cor.~9.5.1);
$\mgeom(\theta_{i_Q^G\sigma})=\mgeom^{\tilde G}(\theta_{\tilde\sigma})$ (Thesis Lemma~7.3.1, \cite{BP} Prop.~4.8.1(i));
conclude by (HYP).
\end{proof}

\begin{proposition}[The defect vanishes on $\Rind$]
\label{defect_eq_zero_of_mem_Rind}
\lean{GSpinGGP.defect_eq_zero_of_mem_Rind}
\leanok
\uses{CentralCharOn, G, GGPData, GSpin, HYP, IsTempered, Rind, UnitaryRep, VirtualTemp, charG, defect, mult, orbitChoice}
Assume (HYP). The defect vanishes on $\Rind(G,\eta)$ (Thesis Prop.~7.4.1(i), extended linearly).
\end{proposition}
\begin{proof}
\leanok
\uses{GammaClass, IsGammaElt, IsQuasiCharacter, LIT-germ_unique, LIT-hc_vMult_eq_mult_parInd, QC, charG_mem_QC, deltaGGP, germ, gspin, jHat, mGeom, mGeom_add, mult_parInd_eq_mGeom, parInd, tendsto_mGeomS, weylDisc}
On a generator: $m(\Pi)=m(i_Q^G\sigma)$, $\mgeom(\Pi)=\mgeom(\theta_{i_Q^G\sigma})$, and these agree.
\end{proof}

\begin{proposition}[Vanishing elliptic character implies vanishing defect]
\label{defect_eq_zero_of_vChar_eq_zero_ell}
\lean{GSpinGGP.defect_eq_zero_of_vChar_eq_zero_ell}
\leanok
\uses{CentralCharOn, G, GGPData, GSpin, HYP, IsTempered, UnitaryRep, VirtualTemp, charG, defect, gspin, mult, orbitChoice}
Assume (HYP). If $\theta_\Pi=0$ on $\ellreg(F)$ then $m(\Pi)=\mgeom(\Pi)$.
\end{proposition}
\begin{proof}
\leanok
\uses{LIT-arthur_mem_Rind_of_vChar_ell, defect_eq_zero_of_mem_Rind}
Such $\Pi$ lie in $\Rind$ (Arthur), where the defect vanishes.
\end{proof}

\begin{theorem}[The spectral expansion, elliptic form]
\label{thm_9_1_1_ell}
\lean{GSpinGGP.thm_9_1_1_ell}
\leanok
\uses{CentralCharOn, G, GGPData, GSpin, HYP, IsTempered, ScuspTestFn, UnitaryRep, VirtualTemp, charG, defect, distribJ, gspin, mGeom, mult, orbitChoice, thetaF}
(Thesis Thm.~10.2.1 = Thm.~1.0.4 and Prop.~7.4.1(ii), elliptic form.) Assume (HYP). For
$f\in\scusp(A_G\backslash G,\eta)$, $J(f)-\mgeom(\theta_f)=\mathrm{defect}(\Pi_f)$, where $\Pi_f$ is a virtual
tempered representation with $\theta_{\Pi_f}=\theta_f$ on $\ellreg(F)$.
\end{theorem}
\begin{proof}
Not yet formalized. The spectral expansion $J(f)=\int_{\mathcal X(G,\chi)}D(\pi)\hat\theta_f(\pi)m(\bar\pi)\,d\pi$
(Thesis Thm.~10.2.1, from Thm.~10.1.1, Arthur's local trace formula, the matricial Paley--Wiener theorem,
Thm.~9.2.1 and Cor.~9.5.1); on the induced part $m=\mgeom$ by Prop.~7.4.1(i).
\end{proof}

\begin{proposition}[The spectral expansion in elliptic form]
\label{exists_virtual_spectral_ell}
\lean{GSpinGGP.exists_virtual_spectral_ell}
\leanok
\uses{CentralCharOn, G, GGPData, GSpin, HYP, IsTempered, ScuspTestFn, UnitaryRep, VirtualTemp, charG, defect, distribJ, gspin, mGeom, mult, orbitChoice, thetaF}
Assume (HYP). For $f\in\scusp(A_G\backslash G,\eta)$ there is a virtual tempered $\Pi_f$ with central character
$\eta$, $\theta_{\Pi_f}=\theta_f$ on $\ellreg(F)$ and $J(f)-\mgeom(\theta_f)=m(\Pi_f)-\mgeom(\Pi_f)$.
\end{proposition}
\begin{proof}
\leanok
\uses{LIT-bp_exists_vChar_eq_thetaF_ell, thm_9_1_1_ell}
Combine the previous theorem with the elliptic virtual character of $\theta_f$.
\end{proof}


% Generated from the Lean dependency data; see the comment at the top of content.tex.
\chapter{\texorpdfstring{The linear form $m_{\mathrm{geom}}$}{The linear form m geom}}
\label{chap:mgeom}
Thesis Prop.~12.1.1(i): convergence, linearity and continuity of $\mgeom$ on quasi-characters.

\begin{lemma}[Measurability of the integrand]
\label{aestronglyMeasurable_torusFn_of_mem_QC}
\lean{GSpinGGP.aestronglyMeasurable_torusFn_of_mem_QC}
\leanok
\uses{CentralCharOn, G, GGPData, GSpin, IsGammaElt, QC, deltaGGP, germ, gspin, orbitChoice, weylDisc}
For $\theta\in\QC(G,\chi)$, the integrand \[ t\mapsto D^G(t)^{1/2}c_{\theta,\mathcal O_t}(t)\Delta(t)^{s-1/2} \] of
$m_{\mathrm{geom},s}$ is a.e.\ strongly measurable on $T(F)/A_{G,H}(F)$ (Thesis \S6.5).
\end{lemma}
\begin{proof}
Not yet formalized. The integrand is continuous on the open set $T_\natural$ (germs are locally constant there) and $0$ off it.
\end{proof}

\begin{lemma}[Boundedness of germs]
\label{germ_bound_of_mem_QC}
\lean{GSpinGGP.germ_bound_of_mem_QC}
\leanok
\uses{CentralCharOn, G, GGPData, GSpin, IsGammaElt, QC, germ, gspin, orbitChoice, weylDisc}
For $\theta\in\QC(G,\chi)$ and a torus $T$, $D^G(t)^{1/2}c_{\theta,\mathcal O_t}(t)$ is bounded on
$T_\natural=\{t\in\Gamma(G,H): T_t=T\}$.
\end{lemma}
\begin{proof}
Not yet formalized. \cite{BP} Prop.~4.5.1 for the individual germs, and compactness of $T(F)/A_{G,H}(F)$.
\end{proof}

\begin{literature}[Integrability of $\Delta^{s-1/2}$]
\label{LIT-bp_integrable_deltaGGP_rpow}
\lean{GSpinGGP.bp_integrable_deltaGGP_rpow}
\leanok
\uses{GGPData, GSpin, IsGammaElt, deltaGGP}
For $x\in\Gamma(G,H)$, $T=T_x$ and $s>0$, $t\mapsto 1_{T_\natural}(t)\Delta(t)^{s-1/2}$ is integrable on
$T(F)/A_{G,H}(F)$.

\assumedfrom{\cite{BP}, Lemma B.1.2(i), applied as in Thesis Prop.~12.1.1(i)}
\end{literature}

\begin{lemma}[Thesis (12.1.2), torus by torus]
\label{isTorusIntegrable_of_mem_QC}
\lean{GSpinGGP.isTorusIntegrable_of_mem_QC}
\leanok
\uses{CentralCharOn, G, GGPData, GSpin, IsGammaElt, QC, deltaGGP, germ, gspin, orbitChoice, weylDisc}
For $\theta\in\QC(G,\chi)$ and $s>0$, the integrand of $m_{\mathrm{geom},s}(\theta)$ is integrable on each
$T(F)/A_{G,H}(F)$.
\end{lemma}
\begin{proof}
\leanok
\uses{LIT-bp_integrable_deltaGGP_rpow, aestronglyMeasurable_torusFn_of_mem_QC, germ_bound_of_mem_QC}
Bounded germs times the integrable $\Delta^{s-1/2}$.
\end{proof}

\begin{lemma}[Existence of the limit]
\label{tendsto_mGeomS_of_mem_QC}
\lean{GSpinGGP.tendsto_mGeomS_of_mem_QC}
\leanok
\uses{CentralCharOn, G, GGPData, GSpin, GammaClass, IsGammaElt, QC, deltaGGP, germ, gspin, orbitChoice, weylDisc}
For $\theta\in\QC(G(F),\chi)$, $\lim_{s\to0^+}m_{\mathrm{geom},s}(\theta)$ exists (Thesis Prop.~12.1.1(i), (12.1.4)).
\end{lemma}
\begin{proof}
Not yet formalized. Induction on $\dim G$: semisimple descent away from $1$, descent to the Lie algebra and homogeneity near $1$.
\end{proof}

\begin{proposition}[$\mgeom$ converges]
\label{tendsto_mGeomS}
\lean{GSpinGGP.tendsto_mGeomS}
\leanok
\uses{CentralCharOn, G, GGPData, GSpin, GammaClass, IsGammaElt, QC, deltaGGP, germ, gspin, orbitChoice, weylDisc}
For $\theta$ in the span $\QC_{\mathrm{fin}}$ of the spaces $\QC(G,\chi)$, the limit $\mgeom(\theta)$ exists
(Thesis Prop.~12.1.1(i)).
\end{proposition}
\begin{proof}
\leanok
\uses{IsQuasiCharacter, LIT-germ_unique, finite_gammaClass, isTorusIntegrable_of_mem_QC, jHat, tendsto_mGeomS_of_mem_QC}
Linearity of $m_{\mathrm{geom},s}$ for $s>0$ and the limit for each $\chi$.
\end{proof}

\begin{proposition}[$\mgeom$ is additive]
\label{mGeom_add}
\lean{GSpinGGP.mGeom_add}
\leanok
\uses{CentralCharOn, G, GGPData, QC, mGeom}
$\mgeom(\theta_1+\theta_2)=\mgeom(\theta_1)+\mgeom(\theta_2)$ on $\QC_{\mathrm{fin}}$ (Thesis Prop.~12.1.1(i)).
\end{proposition}
\begin{proof}
\leanok
\uses{GSpin, GammaClass, IsGammaElt, IsQuasiCharacter, LIT-germ_unique, deltaGGP, finite_gammaClass, germ, gspin, isTorusIntegrable_of_mem_QC, jHat, orbitChoice, tendsto_mGeomS, weylDisc}
From convergence.
\end{proof}

\begin{literature}[Harish-Chandra's compactness lemma]
\label{LIT-hc_exists_invariantClopen_superset}
\lean{GSpinGGP.hc_exists_invariantClopen_superset}
\leanok
\uses{G, GGPData, GSpin}
Every compact $C\subset G(F)$ is contained in an open, closed, invariant, $A_G(F)$-stable $\Omega$ that is compact
modulo conjugation and $A_G(F)$.

\assumedfrom{\cite{HC70}, Part I; \cite{BP}, Prop.~3.1.1(i) (= Thesis Prop.~3.5.1(i)), adapted}
\end{literature}

\begin{proposition}[Support of $\mgeom$]
\label{exists_qcCutoff_mGeom}
\lean{GSpinGGP.exists_qcCutoff_mGeom}
\leanok
\uses{CentralCharOn, G, GGPData, GSpin, QC, gspin, mGeom}
There is an invariant clopen $\Omega$, compact modulo conjugation and $A_G(F)$, such that $\mgeom(\theta)=0$ for
every $\theta\in\QC(G,\chi)$ vanishing on $\Omega\cap\Greg(F)$ (Thesis Prop.~12.1.1(i), continuity).
\end{proposition}
\begin{proof}
\leanok
\uses{GammaClass, IsGammaElt, LIT-germ_unique, LIT-hc_exists_invariantClopen_superset, deltaGGP, finite_gammaClass, germ, jHat, orbitChoice, weylDisc}
$\Gamma(G,H)$ is a finite union of images of tori compact modulo $A_{G,H}(F)$; take $\Omega$ around them; germs only depend on $\theta$ near the point.
\end{proof}


% Generated from the Lean dependency data; see the comment at the top of content.tex.
\chapter{The geometric expansion}
\label{chap:geometric}
Thesis \S12.2--12.5: the geometric expansion $J(f)=\mgeom(\theta_f)$ under (HYP), split into a neighbourhood
of the centre (descent to the Lie algebra) and its complement (semisimple descent).

\section{\texorpdfstring{Extending $J$ to quasi-characters}{Extending J to quasi-characters}}
\begin{definition}[Extensions $J_{\mathrm{qc}}$ of $J$]
\label{IsJqc}
\lean{GSpinGGP.IsJqc}
\leanok
\uses{CentralCharOn, G, GGPData, GSpin, QC, ScuspTestFn, distribJ, gspin, mGeom, thetaF, thetaF_mem_QC}
A linear form $J_{\mathrm{qc}}$ on $\QC(G,\chi^{-1})$ that depends only on values on $\Greg(F)$, satisfies
$J(f)=J_{\mathrm{qc}}(\theta_f)$, is continuous, and $J_{\mathrm{qc}}-\mgeom$ is supported in the elliptic set
(Thesis Prop.~12.2.3). This differs from Thesis Prop.~12.2.3, which has $\Supp(J_{\mathrm{qc}})\subset G(F)_{\mathrm{ell}}$:
here $\Gamma(G,H)$ need not lie in the elliptic set.
\end{definition}

\begin{lemma}[$J-\mgeom$ only sees the elliptic set]
\label{distribJ_sub_mGeom_eq_of_eqOn_ell}
\lean{GSpinGGP.distribJ_sub_mGeom_eq_of_eqOn_ell}
\leanok
\uses{CentralCharOn, G, GGPData, GSpin, HYP, ScuspTestFn, distribJ, gspin, mGeom, thetaF}
Assume (HYP). If $\theta_f=\theta_{f'}$ on $\ellreg(F)$ then $J(f)-\mgeom(\theta_f)=J(f')-\mgeom(\theta_{f'})$
(Thesis Prop.~7.4.1(ii)).
\end{lemma}
\begin{proof}
\leanok
\uses{IsTempered, UnitaryRep, VirtualTemp, charG, defect, defect_eq_zero_of_vChar_eq_zero_ell, exists_virtual_spectral_ell, mult, orbitChoice}
Both are defects of virtual representations with the same elliptic character.
\end{proof}

\begin{proposition}[Existence of $J_{\mathrm{qc}}$]
\label{exists_isJqc}
\lean{GSpinGGP.exists_isJqc}
\leanok
\uses{CentralCharOn, G, GGPData, HYP, IsJqc, QC}
Assume (HYP). Then $J$ extends to some $J_{\mathrm{qc}}$ as above (Thesis Prop.~12.2.3).
\end{proposition}
\begin{proof}
\leanok
\uses{GSpin, GammaClass, H, IsGammaElt, IsQuasiCharacter, LIT-bp_scusp_surj_ell, LIT-bp_thetaF_exists, LIT-germ_unique, LIT-hc_regularSet_open_conull, N, ScuspTestFn, deltaGGP, distribJ, distribJ_add, distribJ_sub_mGeom_eq_of_eqOn_ell, exists_qcCutoff_mGeom, germ, gspin, isThetaF_chiProj, jHat, mGeom, mGeom_add, orbitChoice, tendsto_mGeomS, thetaF, thetaF_mem_QC, thetaF_spec, weylDisc, xi}
$J_{\mathrm{qc}}(\theta)=\mgeom(\theta)+(J(f)-\mgeom(\theta_f))$ for any $f$ with $\theta_f=\theta$ on the elliptic set (elliptic isolation); well defined by the previous lemma.
\end{proof}

\begin{lemma}[$J(f)$ depends only on $\theta_f$]
\label{distribJ_eq_of_thetaF_eqOn}
\lean{GSpinGGP.distribJ_eq_of_thetaF_eqOn}
\leanok
\uses{CentralCharOn, G, GGPData, HYP, ScuspTestFn, distribJ, gspin, thetaF}
Assume (HYP). If $\theta_f=\theta_{f'}$ on $\Greg(F)$ then $J(f)=J(f')$ (Thesis Prop.~12.2.3).
\end{lemma}
\begin{proof}
\leanok
\uses{IsJqc, QC, exists_isJqc, thetaF_mem_QC}
$J=J_{\mathrm{qc}}\circ\theta$.
\end{proof}

\section{Away from the centre}
Thesis \S12.3.

\begin{lemma}[Invariant clopen sets are saturated along unipotent radicals]
\label{mem_mul_parabolicU_iff}
\lean{GSpinGGP.mem_mul_parabolicU_iff}
\leanok
\uses{G, GGPData, GSpin}
For an invariant clopen $\Omega$, $p\in P(F)$ and $u\in U(F)$, $pu\in\Omega\iff p\in\Omega$.
\end{lemma}
\begin{proof}
Not yet formalized. Contract $U$ by a cocharacter of $A_M$; the Levi part lies in the closure of both classes.
\end{proof}

\begin{proposition}[Partition of unity]
\label{geomAt_glue}
\lean{GSpinGGP.geomAt_glue}
\leanok
\uses{CentralCharOn, G, GGPData, GSpin, HYP, ScuspTestFn, distribJ, gspin, mGeom, thetaF}
Assume (HYP). If invariant clopen sets $\Omega_i$ are each either good (the geometric expansion holds for $f$
supported there) or contained in the zero set of $\theta$, and $\theta_f$ is supported in $\bigcup_i\Omega_i$ and
in the support of $\theta$, then $J(f)=\mgeom(\theta_f)$ (Thesis Prop.~12.3.1, \cite{BP} \S11.6).
\end{proposition}
\begin{proof}
\leanok
\uses{GammaClass, H, IsGammaElt, IsQuasiCharacter, LIT-germ_unique, LIT-hc_regularSet_open_conull, N, QC, deltaGGP, distribJ_add, distribJ_eq_of_thetaF_eqOn, germ, isThetaF_chiProj, jHat, mGeom_add, mem_mul_parabolicU_iff, orbitChoice, thetaF_mem_QC, thetaF_spec, weylDisc}
Induction on the number of sets, cutting $f$ by $\Omega_i$.
\end{proof}

\begin{literature}[Neighbourhoods of the near-central set]
\label{LIT-hc_exists_invariantClopen_nearCentral}
\lean{GSpinGGP.hc_exists_invariantClopen_nearCentral}
\leanok
\uses{G, GGPData, GSpin}
For every neighbourhood $U$ of $1$ there is an open, closed, invariant, $A_G(F)$-stable $\Omega$ containing every
near-central element ($1$ in the closure of $A_G(F)c^{G(F)}$), each element of which is conjugate into $U$
modulo $A_G(F)$.

\assumedfrom{\cite{HC70}, Part I; \cite{BP}, \S3.3 (= Thesis \S3.5.2), adapted}
\end{literature}

\begin{lemma}[Thesis Prop.~12.3.1, local form]
\label{exists_good_nhd_of_not_nearCentral}
\lean{GSpinGGP.exists_good_nhd_of_not_nearCentral}
\leanok
\uses{CentralCharOn, G, GGPData, GSpin, HYP, ScuspTestFn, distribJ, gspin, mGeom, thetaF}
Assume (HYP). Every $c$ whose semisimple part is not in $A_G(F)$ has an invariant clopen neighbourhood that
either contains no elliptic regular element or on which the geometric expansion holds.
\end{lemma}
\begin{proof}
Not yet formalized. Semisimple descent to $G_x=G'_x\cdot G''_x$: (HYP) for the smaller GGP triple
$(G'_x,H'_x,\xi'_x)$ and Arthur's local trace formula on $G''_x$ (Thesis (12.3.3), Thm.~4.7.1).
\end{proof}

\begin{proposition}[Thesis Prop.~12.3.1]
\label{geomAt_away}
\lean{GSpinGGP.geomAt_away}
\leanok
\uses{CentralCharOn, G, GGPData, HYP, ScuspTestFn, distribJ, gspin, mGeom, thetaF}
Assume (HYP). If $f\in\scusp(A_G\backslash G,\chi^{-1})$ and $\theta_f$ vanishes near $1$, then $J(f)=\mgeom(\theta_f)$.
\end{proposition}
\begin{proof}
\leanok
\uses{GSpin, H, IsJqc, IsQuasiCharacter, LIT-hc_exists_invariantClopen_nearCentral, LIT-hc_exists_invariantClopen_superset, LIT-hc_regularSet_open_conull, N, QC, exists_good_nhd_of_not_nearCentral, exists_isJqc, geomAt_glue, jHat, thetaF_mem_QC, thetaF_spec}
Cover the support of $\theta_f$ (compact modulo conjugation and $A_G$) by finitely many good sets and glue.
\end{proof}

\section{Near the centre: descent to the Lie algebra}
Thesis \S12.4--12.5 and Ch.~11. Here $\mathfrak g^0=\mathfrak{so}(W)\oplus\mathfrak{so}(V)$.

\begin{definition}[The distribution $J^{\mathrm{Lie}}$]
\label{distribJLie}
\lean{GSpinGGP.GGPData.distribJLie}
\leanok
\uses{G, GGPData, GSpin, H, N, gspin}
$J^{\mathrm{Lie}}(f)=\int_{H(F)\backslash G(F)/A_G(F)}K^{\mathrm{Lie}}(f,x)\,dx$ for $f$ on $\mathfrak g^0(F)$
(Thesis \S10.1; \cite{Luo} \S4, \cite{BP} \S8.1).
\end{definition}

\begin{definition}[$\theta_f$ on the Lie algebra]
\label{thetaFLie}
\lean{GSpinGGP.GGPData.thetaFLie}
\leanok
\uses{G, GGPData, GSpin, distribJLie, gspin}
For $f\in\mathcal S_{\mathrm{scusp}}(\mathfrak g^0(F))$, the quasi-character $\theta_f$ (Thesis \S4.8; \cite{BP} \S5.5).
\end{definition}

\begin{definition}[$m^{\mathrm{Lie}}_{\mathrm{geom}}$]
\label{mGeomLie}
\lean{GSpinGGP.GGPData.mGeomLie}
\leanok
\uses{G, GGPData, GSpin, GammaClass, H, IsGammaElt, N, gspin, jHat, weylDisc}
$m^{\mathrm{Lie}}_{\mathrm{geom}}(\theta)=\lim_{s\to0^+}m^{\mathrm{Lie}}_{\mathrm{geom},s}(\theta)$ (Thesis Prop.~12.1.1(iii)).
\end{definition}

\begin{lemma}[Descent to the Lie algebra near $1$]
\label{exists_lie_descent}
\lean{GSpinGGP.exists_lie_descent}
\leanok
\uses{CentralCharOn, G, GGPData, GSpin, ScuspTestFn, distribJ, distribJLie, gspin, mGeom, mGeomLie, thetaF, thetaFLie}
(Thesis Prop.~12.4.1(i), Prop.~12.1.1(iv).) If neither $W$ nor $V$ is a hyperbolic plane, there is an invariant
clopen $\Omega\ni1$ such that every $f\in\scusp(A_G\backslash G,\chi^{-1})$ supported in $\Omega$ has a strongly cuspidal
$F$ on $\mathfrak g^0$ with $J(f)=J^{\mathrm{Lie}}(F)$ and $\mgeom(\theta_f)=m^{\mathrm{Lie}}_{\mathrm{geom}}(\theta_F)$.
\end{lemma}
\begin{proof}
Not yet formalized. $\Omega=A_G(F)\exp(\omega)$ for a $G(F)$-excellent $\omega$, $F=(j_G^H)^{1/2}f_\omega$.
\end{proof}

\begin{literature}[Spectral expansion of $J^{\mathrm{Lie}}$]
\label{LIT-jLie_eq_of_thetaFLie_eqOn}
\lean{GSpinGGP.jLie_eq_of_thetaFLie_eqOn}
\leanok
\uses{G, GGPData, GSpin, distribJLie, gspin}
(Thesis Thm.~11.6.1, in normalization-free form.) If neither $W$ nor $V$ is a hyperbolic plane, $J^{\mathrm{Lie}}(f)$
depends only on $\theta_f|_{\mathfrak g^0_{\mathrm{reg}}}$.

\assumedfrom{\cite{Luo}, Thm.~6.6.1; \cite{Wal10}, \S9--10; following \cite{BP}, Thm.~10.8.1}
\end{literature}

\begin{lemma}[$J^{\mathrm{Lie}}$ factors through $\theta_f$]
\label{distribJLie_eq_jqcLie}
\lean{GSpinGGP.distribJLie_eq_jqcLie}
\leanok
\uses{G, GGPData, GSpin, distribJLie, gspin, thetaFLie}
If neither $W$ nor $V$ is a hyperbolic plane, $J^{\mathrm{Lie}}(f)=J^{\mathrm{Lie}}_{\mathrm{qc}}(\theta_f)$ for every
strongly cuspidal $f$ on $\mathfrak g^0$ (Thesis Prop.~12.4.1(ii)).
\end{lemma}
\begin{proof}
\leanok
\uses{LIT-jLie_eq_of_thetaFLie_eqOn}
From the spectral expansion of $J^{\mathrm{Lie}}$.
\end{proof}

\begin{lemma}[Thesis Prop.~12.5.1, test-function form]
\label{prop_11_8_1_of_jqc}
\lean{GSpinGGP.prop_11_8_1_of_jqc}
\leanok
\uses{G, GGPData, GSpin, H, HYP, N, distribJLie, gspin, jHat, mGeomLie, thetaFLie}
Assume (HYP), and that $J^{\mathrm{Lie}}$ factors through $\theta_f$ by a form seeing only $\mathfrak g^0_{\mathrm{reg}}$.
Then $J^{\mathrm{Lie}}(f)-m^{\mathrm{Lie}}_{\mathrm{geom}}(\theta_f)=\sum_{\mathcal O\in\Nilreg(\mathfrak g^0)}c_{\mathcal O}\,c_{\theta_f,\mathcal O}(0)$
for some constants $c_{\mathcal O}$.
\end{lemma}
\begin{proof}
Not yet formalized. The difference is a homogeneous continuous form, vanishing on quasi-characters with
$0\notin\Supp\theta$ by semisimple descent and (HYP) (Prop.~12.4.1(iv)); such a form is a combination of germs at
$0$ (\cite{BP} Prop.~4.6.1).
\end{proof}

\begin{proposition}[Thesis Prop.~12.5.1]
\label{prop_11_8_1}
\lean{GSpinGGP.prop_11_8_1}
\leanok
\uses{G, GGPData, GSpin, H, HYP, N, distribJLie, gspin, jHat, mGeomLie, thetaFLie}
Assume (HYP), neither $W$ nor $V$ a hyperbolic plane. Then
$J^{\mathrm{Lie}}(f)-m^{\mathrm{Lie}}_{\mathrm{geom}}(\theta_f)=\sum_{\mathcal O\in\Nilreg(\mathfrak g^0)}c_{\mathcal O}\,c_{\theta_f,\mathcal O}(0)$.
\end{proposition}
\begin{proof}
\leanok
\uses{distribJLie_eq_jqcLie, prop_11_8_1_of_jqc}
The test-function form with $J^{\mathrm{Lie}}_{\mathrm{qc}}$.
\end{proof}

\begin{literature}[The endgame]
\label{LIT-lie_endgame}
\lean{GSpinGGP.lie_endgame}
\leanok
\uses{G, GGPData, GSpin, H, N, distribJLie, gspin, jHat, mGeomLie, thetaFLie}
(Thesis \S12.5.) If $J^{\mathrm{Lie}}(f)-m^{\mathrm{Lie}}_{\mathrm{geom}}(\theta_f)=\sum_{\mathcal O}c_{\mathcal O}c_{\theta_f,\mathcal O}(0)$
for every strongly cuspidal $f$ on $\mathfrak g^0$, then every $c_{\mathcal O}$ vanishes.

\assumedfrom{\cite{Wal10}, \S11.4--11.6; \cite{Luo}, Prop.~7.3.7, Lemma B.3.1}
\end{literature}

\begin{theorem}[$J^{\mathrm{Lie}}=J^{\mathrm{Lie}}_{\mathrm{geom}}$]
\label{distribJLie_eq_mGeomLie}
\lean{GSpinGGP.distribJLie_eq_mGeomLie}
\leanok
\uses{G, GGPData, GSpin, HYP, distribJLie, gspin, mGeomLie, thetaFLie}
Assume (HYP), neither $W$ nor $V$ a hyperbolic plane. Then $J^{\mathrm{Lie}}(f)=m^{\mathrm{Lie}}_{\mathrm{geom}}(\theta_f)$
for strongly cuspidal $f$ on $\mathfrak g^0$ (Thesis Thm.~12.2.2).
\end{theorem}
\begin{proof}
\leanok
\uses{H, LIT-lie_endgame, N, jHat, prop_11_8_1}
Prop.~12.5.1 and the endgame.
\end{proof}

\begin{lemma}[The case $V$ a hyperbolic plane]
\label{geomAt_near_of_isHypPlane_V}
\lean{GSpinGGP.geomAt_near_of_isHypPlane_V}
\leanok
\uses{CentralCharOn, G, GGPData, GSpin, ScuspTestFn, distribJ, mGeom, thetaF}
If $V$ is a hyperbolic plane ($\dim W=1$, $r=0$), the geometric expansion holds on all of $G(F)$.
\end{lemma}
\begin{proof}
\leanok
\uses{GammaClass, H, IsGammaElt, LIT-exists_isSelfDual, N, deltaGGP, germ, gspin, jHat, orbitChoice, weylDisc, xi}
$G$ is a split torus, $A_G=G$; then $J(f)=f(1)$, $\theta_f=f$ and $\mgeom(\theta_f)=c_{\theta_f,0}(1)=f(1)$, using a self-dual measure for $\hat j(0,\cdot)=1$.
\end{proof}

\begin{lemma}[The case $W$ a hyperbolic plane]
\label{geomAt_near_of_isHypPlane_W}
\lean{GSpinGGP.geomAt_near_of_isHypPlane_W}
\leanok
\uses{CentralCharOn, G, GGPData, GSpin, HYP, ScuspTestFn, distribJ, mGeom, thetaF}
Assume (HYP) and $W$ a hyperbolic plane. There is an invariant clopen $\Omega\ni1$ on which the geometric
expansion holds.
\end{lemma}
\begin{proof}
Not yet formalized. Here $A_{\bar G}=\SO(W)$ is a nontrivial split central torus; the Lie-algebra argument must be
run modulo $A_{\bar G}$, which the cited literature does not state in this form (\cite{BP} Ch.~10--11 modulo $A_{\bar G}$).
\end{proof}

\begin{theorem}[The geometric expansion near the centre]
\label{geomAt_near}
\lean{GSpinGGP.geomAt_near}
\leanok
\uses{CentralCharOn, G, GGPData, GSpin, HYP, ScuspTestFn, distribJ, mGeom, thetaF}
(Thesis Prop.~12.4.1 and Thm.~12.2.2.) Assume (HYP). For every $\chi$ there is an invariant clopen $\Omega\ni1$
with $J(f)=\mgeom(\theta_f)$ for every $f\in\scusp(A_G\backslash G,\chi^{-1})$ supported in $\Omega$.
\end{theorem}
\begin{proof}
\leanok
\uses{distribJLie, distribJLie_eq_mGeomLie, exists_lie_descent, geomAt_near_of_isHypPlane_V, geomAt_near_of_isHypPlane_W, gspin, mGeomLie, thetaFLie}
Cases: $V$ a hyperbolic plane; $W$ a hyperbolic plane; otherwise descend to $\mathfrak g^0$ and use $J^{\mathrm{Lie}}=J^{\mathrm{Lie}}_{\mathrm{geom}}$.
\end{proof}


% Generated from the Lean dependency data; see the comment at the top of content.tex.
\chapter{Further statements not yet used by the main theorem}
\label{chap:offpath}
These statements are formalized but not (yet) on the path of the main theorem: the analytic estimates of
Thesis Ch.~8 and 10 (the convergence of $J$ is stated separately above), the Levi decomposition and transport
of $\mgeom$ (Thesis \S7.2, \S9.3), and checks of the normalization of $\hat j$ on $\mathfrak{sl}_2$.

\section{Estimates}
Thesis Ch.~8, \S10.1.

\begin{proposition}[Good minimal parabolic subgroups]
\label{exists_good_minimal_parabolic}
\lean{GSpinGGP.GGPData.exists_good_minimal_parabolic}
\leanok
\uses{G, GGPData, GSpin, H}
(Thesis Prop.~8.1.1(i), existence.) There are full isotropic flags of $W$ and $V$ whose stabilizer $\bar P$
satisfies $H\cap\bar U_{\min}=\{1\}$, with $H(F)\bar P(F)$ open in $G(F)$ and of Haar-null complement.
\end{proposition}
\begin{proof}
Not yet formalized (adapted from \cite{BP} Prop.~6.4.1). ``Good'' is replaced by its two consequences used in the estimates.
\end{proof}

\begin{lemma}[Thesis Lemma~8.2.1(i)]
\label{lemma_6_5_1_i}
\lean{GSpinGGP.GGPData.lemma_6_5_1_i}
\leanok
\uses{G, GGPData, GSpin, H, N}
For $r=0$ there is $\epsilon>0$ with $\int_{H_0(F)/A_{G,H}(F)}\Xi^{G_0}(h_0)e^{\epsilon\sigma(h_0)}\,dh_0<\infty$.
This differs from Thesis Lemma~8.2.1(i), which integrates over $H_0(F)/A_{H_0}(F)$.
\end{lemma}
\begin{proof}
Not yet formalized (adapted from \cite{BP} Lemma 6.5.1(i)).
\end{proof}

\begin{lemma}[Thesis Lemma~8.2.1(ii)]
\label{lemma_6_5_1_ii}
\lean{GSpinGGP.GGPData.lemma_6_5_1_ii}
\leanok
\uses{G, GGPData, GSpin, H, N}
There is $d>0$ with $\int_{H(F)/A_{G,H}(F)}\Xi^G(h)\sigma(h)^{-d}\,dh<\infty$.
\end{lemma}
\begin{proof}
Not yet formalized (adapted from \cite{BP} Lemma 6.5.1(ii)).
\end{proof}

\begin{lemma}[Thesis Lemma~8.2.1(iii)]
\label{lemma_6_5_1_iii}
\lean{GSpinGGP.GGPData.lemma_6_5_1_iii}
\leanok
\uses{G, GGPData, GSpin, H, N}
For every $\delta>0$ there is $\epsilon>0$ with
$\int_{H(F)/A_{G,H}(F)}\Xi^G(h)e^{\epsilon\sigma(h)}(1+|\lambda(h)|)^{-\delta}\,dh<\infty$.
\end{lemma}
\begin{proof}
Not yet formalized (adapted from \cite{BP} Lemma 6.5.1(iii)).
\end{proof}

\begin{theorem}[Thesis Thm.~10.1.1(i)]
\label{kernelK_bound}
\lean{GSpinGGP.GGPData.kernelK_bound}
\leanok
\uses{CentralCharOn, G, GGPData, GSpin, H, N, xi}
There are $d>0$ and a continuous seminorm $\nu$ on $\mathcal C(A_G(F)\backslash G(F),\chi^{-1})$ with
$|K(f,x)|\le\nu(f)\,\Xi^{H\backslash G}(x)^2\sigma_{H\backslash G}(x)^{d}$.
\end{theorem}
\begin{proof}
Not yet formalized (adapted from \cite{BP} Thm.~8.1.1(i)).
\end{proof}

\begin{theorem}[Thesis Thm.~10.1.1(ii)]
\label{kernelK_bound_scusp}
\lean{GSpinGGP.GGPData.kernelK_bound_scusp}
\leanok
\uses{CentralCharOn, G, GGPData, GSpin, H, N, xi}
For every $d>0$ there is a continuous seminorm $\nu_d$ with
$|K(f,x)|\le\nu_d(f)\,\Xi^{H\backslash G}(x)^2\sigma_{H\backslash G}(x)^{-d}$ for strongly cuspidal $f$.
\end{theorem}
\begin{proof}
Not yet formalized (adapted from \cite{BP} Thm.~8.1.1(ii)).
\end{proof}

\section{Levi subgroups}
\begin{proposition}[The Levi decomposition]
\label{exists_leviHom_of_leviSpec}
\lean{GSpinGGP.GGPData.exists_leviHom_of_leviSpec}
\leanok
\uses{G, GGPData, GSpin}
For a GGP datum $\tilde d$ realizing the Levi pair $(\tilde W,\tilde V)$ of a pair of isotropic flags, there is a
surjective homomorphism with central kernel from the Levi quotient $P/U$ onto
$\prod\mathrm{GL}(Z'_{i,W})\times\prod\mathrm{GL}(Z'_{j,V})\times\tilde G(F)$ (Thesis \S5.1.2, (9.3.1)--(9.3.3)).
\end{proposition}
\begin{proof}
Not yet formalized: it needs the Clifford-algebra description of the Levi subgroup.
\end{proof}

\begin{lemma}[Transport of $\mgeom$]
\label{mGeom_congr_iso}
\lean{GSpinGGP.GGPIso.mGeom_congr_iso}
\leanok
\uses{G, GGPData, GSpin, H, IsQuasiCharacter, gspin, mGeom}
For an isomorphism of GGP triples $e:G\to G'$ and a quasi-character $\theta'$ of $G'(F)$,
$\mgeom(\theta'\circ e)=\mgeom(\theta')$ (Thesis \S7.2, Prop.~12.1.1).
\end{lemma}
\begin{proof}
Not yet formalized: needs the transport of $\mathcal O_x$ and $\hat j$.
\end{proof}

\section{\texorpdfstring{Normalization checks on $\mathfrak{sl}_2$}{Normalization checks on sl 2}}
Sanity checks of the definition of $\hat j$ for $\mathrm{GL}_2(F)$
acting on $\mathfrak{sl}_2(F)$ with the trace form.

\begin{lemma}[The regular orbit measure exists]
\label{exists_isOrbitMeasure_regOrbit}
\lean{GSpinGGP.Sanity.NilpotentOrbit.exists_isOrbitMeasure_regOrbit}
\leanok
The $\mathrm{GL}_2(F)$-invariant measure on the regular nilpotent orbit, normalized by the self-dual measure on
$\mathfrak g/\mathfrak g_E$ (\cite{BP} \S1.6, Thesis \S2.4), exists.
\end{lemma}
\begin{proof}
Not yet formalized: needs the $p$-adic inverse function theorem and invariant measures on homogeneous spaces.
\end{proof}

\begin{lemma}[$\hat j(\mathcal O_{\mathrm{reg}},\cdot)$ on the split torus]
\label{jHat_regOrbit_split}
\lean{GSpinGGP.Sanity.NilpotentOrbit.jHat_regOrbit_split}
\leanok
\uses{jHat}
For $X=x\,\mathrm{diag}(1,-1)$, $x\neq0$: $\hat j(\mathcal O_{\mathrm{reg}},X)=2/|2x|$.
\end{lemma}
\begin{proof}
Not yet formalized (hand computation, every residue characteristic).
\end{proof}

\begin{lemma}[$\hat j(\mathcal O_{\mathrm{reg}},\cdot)$ on elliptic tori]
\label{jHat_regOrbit_elliptic}
\lean{GSpinGGP.Sanity.NilpotentOrbit.jHat_regOrbit_elliptic}
\leanok
\uses{jHat}
For $\varepsilon$ a nonsquare and $x\neq0$, $\hat j(\mathcal O_{\mathrm{reg}},x\begin{pmatrix}0&1\\\varepsilon&0\end{pmatrix})=0$.
\end{lemma}
\begin{proof}
Not yet formalized: $\mu_{\mathcal O_{\mathrm{reg}}}$ is induced from the Borel, so its Fourier transform is induced from the split torus.
\end{proof}



\begin{thebibliography}{99}
\bibitem{Thesis} S.~Pease, \emph{The Local Gan--Gross--Prasad Conjecture for General Spin Groups},
  PhD dissertation, Rutgers University--Newark, 2026, doi:10.7282/t3-qb3s-s249.
\bibitem{Art91} J.~Arthur, \emph{A local trace formula}, Publ.\ Math.\ IH\'ES 73 (1991), 5--96.
\bibitem{Art93} J.~Arthur, \emph{On elliptic tempered characters}, Acta Math.\ 171 (1993), 73--138.
\bibitem{Artin} E.~Artin, \emph{Geometric Algebra}, Interscience, 1957.
\bibitem{BP} R.~Beuzart-Plessis, \emph{A local trace formula for the Gan--Gross--Prasad conjecture for
  unitary groups: the archimedean case}, Ast\'erisque 418, 2020.
\bibitem{Dixmier} J.~Dixmier, \emph{$C^*$-algebras}, North-Holland, 1977.
\bibitem{ET} M.~Emory, S.~Takeda, \emph{Contragredients and a multiplicity one theorem for general
  spin groups}.
\bibitem{Folland} G.~Folland, \emph{A Course in Abstract Harmonic Analysis}, CRC Press.
\bibitem{GGP} W.~T.~Gan, B.~Gross, D.~Prasad, \emph{Symplectic local root numbers, central critical
  $L$-values, and restriction problems in the representation theory of classical groups},
  Ast\'erisque 346, 2012.
\bibitem{HC70} Harish-Chandra, \emph{Harmonic analysis on reductive $p$-adic groups}, LNM 162,
  Springer, 1970.
\bibitem{HC99} Harish-Chandra, \emph{Admissible invariant distributions on reductive $p$-adic groups},
  University Lecture Series 16, AMS, 1999.
\bibitem{Knapp} A.~Knapp, \emph{Lie Groups Beyond an Introduction}, 2nd ed., Birkh\"auser, 2002.
\bibitem{Kot83} R.~Kottwitz, \emph{Sign changes in harmonic analysis on reductive groups},
  Trans.\ Amer.\ Math.\ Soc.\ 278 (1983), 289--297.
\bibitem{Lam} T.~Y.~Lam, \emph{Introduction to Quadratic Forms over Fields}, AMS, 2005.
\bibitem{Luo} Z.~Luo, \emph{A local trace formula for the local Gan--Gross--Prasad conjecture for
  special orthogonal groups}, arXiv:2009.13947.
\bibitem{MW87} C.~M\oe glin, J.-L.~Waldspurger, \emph{Mod\`eles de Whittaker d\'eg\'en\'er\'es pour
  des groupes $p$-adiques}, Math.\ Z.\ 196 (1987), 427--452.
\bibitem{Rodier} F.~Rodier, \emph{Mod\`ele de Whittaker et caract\`eres de repr\'esentations},
  LNM 466, Springer, 1975.
\bibitem{SerreCG} J.-P.~Serre, \emph{Cohomologie galoisienne}, LNM 5, Springer.
\bibitem{SerreLF} J.-P.~Serre, \emph{Local Fields}, GTM 67, Springer, 1979.
\bibitem{Steinberg} R.~Steinberg, \emph{Endomorphisms of linear algebraic groups}, Mem.\ Amer.\ Math.\
  Soc.\ 80, 1968.
\bibitem{vD72} G.~van Dijk, \emph{Computation of certain induced characters of $p$-adic groups},
  Math.\ Ann.\ 199 (1972), 229--240.
\bibitem{Wal01} J.-L.~Waldspurger, \emph{Int\'egrales orbitales nilpotentes et endoscopie pour les
  groupes classiques non ramifi\'es}, Ast\'erisque 269, 2001.
\bibitem{WalPl} J.-L.~Waldspurger, \emph{La formule de Plancherel pour les groupes $p$-adiques
  d'apr\`es Harish-Chandra}, J.\ Inst.\ Math.\ Jussieu 2 (2003), 235--333.
\bibitem{Wal10} J.-L.~Waldspurger, \emph{Une formule int\'egrale reli\'ee \`a la conjecture locale de
  Gross--Prasad}, Compositio Math.\ 146 (2010), 1180--1290.
\bibitem{Weil} A.~Weil, \emph{Basic Number Theory}, Springer.
\bibitem{Yan} P.~Yan, \emph{Uniqueness of Bessel models for GSpin}, arXiv:2503.20116.
\end{thebibliography}
